% =====================================================================
%  Harald Høffding, Philosophy in Sweden
%  Translated from the Danish: „Filosofien i Sverig“, Nordisk tidskrift för vetenskap, konst och industri, [Ny serie] Andra årgången (Stockholm 1879), pp. 30--73.
%
%  TRANSLATION -- GENERATED; DO NOT EDIT BY HAND.  Built by ebook-method/hoeffding-1879-82/tr_parts.py build filosofien-i-sverig-1879 --write
%  from the parts in ebook-method/hoeffding-1879-82/tr/ (GLOSSARY.md, TRANSLATOR-BRIEF.md).  Made from the collated
%  transcription texts/hoeffding/filosofien-i-sverig-1879/transcription.tex (2026-10-15); [N] in the margin = the page of the original.
%  Terms: GLOSSARY.md (choices marked ◆ there for Hans to review).
% =====================================================================
\documentclass[12pt,a4paper,openany]{book}
\usepackage[utf8]{inputenc}
\usepackage[T1]{fontenc}
\usepackage{libertinus}
\usepackage{libertinust1math}
\usepackage[german,french,english]{babel}
\usepackage{textalpha}
\usepackage[protrusion=true,expansion=true]{microtype}
\usepackage{setspace}
\usepackage[margin=1.3in]{geometry}
\usepackage{fancyhdr}
\usepackage{hyperref}
\hypersetup{pdftitle={Philosophy in Sweden}, pdfauthor={Harald Høffding}, hidelinks}
\pagestyle{fancy}
\fancyhf{}
\fancyhead[LE,RO]{\thepage}
\fancyhead[C]{\textit{H. Høffding: Philosophy in Sweden (1879)}}
\fancypagestyle{plain}{\fancyhf{}\fancyhead[LE,RO]{\thepage}}
\setstretch{1.3}
\usepackage{marginnote}
\usepackage{graphicx}
\DeclareUnicodeCharacter{0254}{\reflectbox{c}}   % ɔ (id est)
\renewcommand{\thefootnote}{\ifcase\value{footnote}\or *\or **\or ***\else\arabic{footnote}\fi)}   % the print marks each page\'s notes *), **), ***)
\newcommand{\opage}[1]{\marginnote{\footnotesize\textit{[#1]}}}
\newcommand{\signature}[1]{\par\smallskip\hfill #1\hspace*{2em}\par}
\pdfstringdefDisableCommands{\renewcommand{\footnote}[2][]{}\renewcommand{\opage}[1]{}}

\begin{document}

\begin{titlepage}
\begin{center}
{\Huge\bfseries Nordisk tidskrift}\\[0.4em]{\large för vetenskap, konst och industri}\\[0.4em]{\large [Nordic Journal of Science, Art and Industry]}\\[0.6em]{\small published by the Letterstedt Society}\\[0.4em]{\small edited by Claes Annerstedt (Uppsala), C. M. Guldberg (Kristiania) and Julius Lange (Copenhagen)}\\[0.6em]1879\\[0.4em]{\small Stockholm: P. A. Norstedt \& Söner}\\[2em]
{\small Translated from the Danish}
\end{center}
\end{titlepage}
\clearpage

% --- p. 30 ---
\setcounter{footnote}{0}%
\opage{30}\begin{center}{\Large\sffamily\bfseries PHILOSOPHY IN SWEDEN.}\\[0.4em]{\bfseries {\small BY} HARALD HÖFFDING.}\\[0.4em]\rule{3em}{0.4pt}\end{center}\medskip

\begin{small}
1. \textit{Axel Nyblæus:} Den filosofiska forskningen i Sverige från slutet af adertonde
århundradet, framställd i sitt sammanhang med filosofiens allmänna utveckling
[Philosophical Research in Sweden from the End of the Eighteenth Century, Presented in Its Connection with the General Development of Philosophy]. Första
delen [First Part]. I. Lund 1873. --- II. Lund 1875.

2. \textit{Axel Nyblæus:} Trenne filosofiska uppsatser [Three Philosophical Essays]. Ny upplaga [New edition]. Lund 1878.

3. \textit{Samuel Grubbe:} Filosofiska skrifter i urval [Selected Philosophical Writings]. Utgifna af [Edited by] \textit{Axel Nyblæus.} I—II.
Praktiska filosofiens historia [History of Practical Philosophy]. Lund 1876—77.

4. \textit{Christopher Jacob Boström:} Om religionens, vishetens och dygdens begrepp
samt deras förhållande till hvarandra [On the Concepts of Religion, Wisdom, and Virtue and Their Relation to One Another]. Öfversättning [from the Latin] af [translation by] \textit{L. Bygdén.}
Upsala 1874.

5. \textit{Waldemar Dons:} Om Boströmianismen [On Boströmianism]. Kristiania 1874.

6. \textit{E---f---t:} Granskning af kandidaten Waldemar Dons' framställning af den
Boströmska filosofien [Examination of Candidate Waldemar Dons's Presentation of the Boströmian Philosophy]. Upsala 1875.
\end{small}
\begin{center}\rule{3em}{0.4pt}\end{center}

Attempts have been made to present the history of philosophy as a continuous
chain of links of thought, each of which engages with the next. Indeed, Hegel even
thought he could demonstrate that the historical development of thinking corresponds
closely to the places of the concepts in the logical system and followed
the latter's laws. Even if one regards this as an arbitrary hypothesis, which can be carried
through only by doing violence to real history, --- even if one has realized
that very different laws hold in logic and in history, one is nonetheless still
inclined to conceive of the development of philosophy in such a way that one philosopher
takes up where the other has left off. The greatest philosopher then becomes,
as Stuart Mill has said, the one who has best studied his predecessor.
It cannot now be denied that the history of philosophy offers examples
enough of such a steadily progressing development, in particular within
each of the great main schools. When a rich watchword
has been given, when a new path has been broken, thinking has for the time being enough
to do in running out the line, exhausting the possibilities that the new
school offers. But what of the thinkers themselves who give the watchword and
inaugurate the new periods? They have had no predecessors to study
and must therefore have received their impulses from sources that lie outside
the historically developed philosophy. It is precisely what is peculiar to
genius, that it sees new possibilities and does not content itself with
drawing conclusions from handed-down principles.

% --- p. 31 ---
\setcounter{footnote}{0}%
\opage{31}Since there are thus places where the historical chain is broken, we are compelled
to acknowledge that philosophy itself cannot be an entirely
independent endeavor, complete in itself. It does not have all its conditions of development
in itself, but points beyond itself, in particular
at the decisive turning points. What the external conditions are that
intervene as determining factors, Plato has in a way already indicated when
he (in the ``Phaedrus'') presents Urania and Calliope as the muses
of philosophy. It is nature and history that give thinking its content
and its impulses. Far from being autonomous, from being able to draw
everything from itself, philosophy therefore presupposes the other sciences,
especially natural science, and it stands at the same time in the most intimate connection
with the ethical and cultural-historical endeavors and
tasks of human life. The turning points in the history of philosophy will prove to be conditioned
by the taking up of new results and methods from natural science,
or by the ethical and cultural-historical striving giving thinking new
tasks. In Socrates it was the ethical interest, the endeavor
after self-knowledge, after clear insight into ``the beautiful and good,'' that set
thought in motion; and that he became the founder of logic is closely
connected with his having become the founder of ethics. Descartes
sought to introduce the mathematical method into philosophy and to develop
a mechanical conception of nature. Locke applied induction and analysis
to the investigation of the origin and development of human knowledge,
and he worked in the field of practical philosophy in
the service of political freedom and toleration. Kant was strongly influenced from his youth
by Newton, and sought in the latter's spirit to present
knowledge through experience as the only valid knowledge and to demonstrate its foundation
in the original organization of consciousness; and in the results he
reached he found at the same time satisfaction for the rightly understood
demands of the ethical ideal.

But Calliope and Urania do not always get along well; they often make
conflicting demands and go each her own way, and Calliope in particular has
often claimed a privileged position and refused to share terms
with her sister. Idealism, which has its source in the ethical
endeavors and tasks of human life and in the feeling of their value, and naturalism,
which proceeds from given reality, which it seeks to penetrate by observation and
analysis, seem to be irreconcilable opposites.
Idealism will perhaps not even leave external nature
to its opponent; it believes itself able to do justice there as well
(Plato, Schelling, Hegel), or it turns away from all study of nature
(Socrates, Fichte). But in any case it reserves for itself
the spiritual domain, and regards it as an encroachment when one wants here to
apply methods similar to those in the domain of external nature. This
opposition runs through the whole history of philosophy and is surely
the deepest ground of all the struggles of which the latter tells. The
% --- p. 32 ---
\setcounter{footnote}{0}%
\opage{32}ideal striving, the feeling of the value of life, especially of the value of spiritual life, we
will not and cannot give up. But on the other hand, observing
and analyzing inquiry makes equally definite demands; it has extended
its conquests to all territories in the domain of external nature,
and now demands with ever greater authority also to gain access to
the world of spirit --- or rather: it has long since, without
asking leave, obtained this access. In its impatience it may
have committed many encroachments, hastily applied methods that
were not sufficiently adapted to the peculiarity of the new domain,
and relied on results that did not have the certainty and range
that were attributed to them. But the goal at which it aims cannot be given up:
the essential unity and coherence of our knowledge, both as regards
method and results. Idealism has been no less guilty of encroachments, and it has in particular confused the ideal
values themselves with the determinate forms, ideas and symbols in which they have at certain
times clothed themselves. It will cease to regard naturalism as
its enemy only when it has given up all forms, so that
they do not become barriers to inquiring thought, and when it has realized
that the true, great and lofty does not lose its luster because it, like
everything else in existence, presupposes determinate conditions and develops
according to determinate laws.

Only rarely have the two tendencies worked in a philosopher with
such harmony as in Kant; naturalism and idealism interpenetrate
in him in a peculiar way. This became possible
because he sought down to the common source of both tendencies: human
nature. So long as they are allowed to unfold freely each by itself,
the opposition increases; the more one approaches the point of departure, the
more they converge. It was precisely Kant's brilliant discovery that all
our thoughts and feelings are determined and borne by the original
organization of our nature, by our standpoint in the universe, and that one must therefore, in order to
understand the validity of our conception of the world, our knowledge and our
valuation of things, descend to this foundation and take it
into consideration. Hereby it becomes possible to summon the contending
world-views and methods before a tribunal that they must all
acknowledge. In a certain sense philosophy has therefore first been constituted
by Kant. No one has, as he did, pointed to its innermost vital
nerve.

That the contending tendencies have asserted themselves again after Kant's time
cannot surprise us, when one rightly considers partly the greatness of
the task he set, partly human impatience, and
partly also that he himself did not express himself with complete clarity, and that
he offered sides from which he could, not without right, be claimed
by one or another of the contending parties. Harmony between the two tendencies
is the great goal toward which we steer; but on the way
% --- p. 33 ---
\setcounter{footnote}{0}%
\opage{33}there will be room and need for many hypotheses and systems.
Every capable and original attempt to carry through a world-view
from one or the other point of departure will, directly or
indirectly, positively or negatively, contribute to progress.

The \textit{Scandinavian countries} have taken only slight independent part in
this great discussion. The waves from the great countries of culture have here,
as in other areas, propagated themselves to us and have called forth lively
interest and capable assimilation, but we have, just as little in the philosophical
as in the religious domain, produced any genius of the first
rank. Yet Sweden surely stands here above the two brother countries. Neither of
these can show such a series of energetic and independent thinkers
as Sweden possesses in Thorild, Höijer, Geijer and Boström, to
name only the most important. The philosophical interest --- which could not
be homeless in a country that had produced spirits such as Saint
Birgitta and Emanuel Swedenborg --- has even in Sweden led
to the development of a peculiar system: the Boströmian, which is now
the ruling one at the Swedish universities, and in which the philosophical
movement seems to have found its provisional resting place, just as
in its time German speculation found its in Hegel's system.

Swedish philosophy has been little known outside its homeland.
As regards Boström in particular, he worked predominantly by
the spoken word; he states expressly in a polemical pamphlet that his
task is to be ``a living teacher for academic youth,
not a writing teacher for the Swedish public''; for the latter
he lacked ``time, health, and ability,'' and it was enough for him that his
philosophy was judged and criticized by his disciples and hearers.\footnote[1]{Mr. \textit{Dons} holds that this predominantly oral tradition has produced a certain faith in authority among Boström's adherents. Mr. \textit{E---f---t} defends the ``living word'' against this charge (pag. 83):
% sv: Tvärtom har det muntliga och lefvande ordet mer lifskraft och är mer verksamt än det döda bokstafsordet. Och detta gäller isynnerhet och i egentlig mening, då ordet har andligt, moraliskt och religiöst, öfver hufvud taget filosofiskt innehåll. Men att det lefvande ordet är mer verksamt än det döda bokstafsordet, måste betyda, att det verkar mer renande och förädlande på sinnet, mer förklarande och upplysande på förståndet, mer stärkande och solliciterande på viljan, för hvilken ett sådant ord är en kraft till lif och det goda. Den verkan, som utgår från ett sådant ord, är sålunda ett högre och klarare sjelfmedvetande.
``On the contrary, the spoken and living word has more vital force and is more effective than the dead word of the letter. And this holds especially and in the proper sense when the word has spiritual, moral and religious, in general philosophical content. But that the living word is more effective than the dead word of the letter must mean that it works more purifyingly and ennoblingly on the mind, more explanatorily and enlighteningly on the understanding, more strengthening and urging on the will, for which such a word is a power toward life and the good. The effect that proceeds from such a word is thus a higher and clearer self-consciousness.'' --- True and beautiful as this is, and though Mr. E---f---t can here appeal to Plato, it is nevertheless surely the case that oral tradition is not favorable to independent criticism. The hearer will not so easily be able to free himself from the personality of the teacher and the ethical and aesthetic impression it makes, whereas the reader can abstract from the personal relation and examine and judge the matter itself more impartially.} His
writings were few and little circulated. In Sweden itself Boström's doctrine has
been subjected to an instructive critique by Prof. \textit{Borelius} in Lund, essentially
from a Hegelian standpoint. This critique called forth an exceedingly
sharp and mocking reply from Boström himself. If the old saying holds true:
``You are wrong, Zeus, for you grow angry!'' --- then the critique must have
hit its mark. --- Nor has there been any lack of attempts from the neighboring countries
to become acquainted with his doctrine. A young Danish
% --- p. 34 ---
\setcounter{footnote}{0}%
\opage{34}scholar, cand. \textit{Kristian Møller,} traveled in the year 1869, at the urging of Orla Lehmann,
to Upsala to study Boströmianism; but a far too
early death prevented him from presenting the result of this journey.
Later Mr. Dons of Kristiania has published a presentation of Boströmianism
and criticized it from a Hegelian standpoint. Against
him Mr. E---f---t has, in a beautifully written treatise, defended the
Boströmian doctrine. The dispute between these two authors concerns (like
the dispute between Boström and Borelius) essentially the question whether
Boströmianism represents an advance or a regression in relation
to Hegelianism. In what follows we shall come to this question
in examining the relation of Swedish speculation to
German.

Of particular importance for a more general and more thorough knowledge
of Swedish philosophy is the fact that Prof. Axel Nyblæus in Lund has
begun to write its history. Prof. Nyblæus, who belongs to the
most significant investigators of the Boströmian school, is known as the author
of several treatises in the philosophy of right and the philosophy of religion. He
has, among other things, in his time contributed an exceedingly interesting and significant
piece to the dispute over faith and knowledge that was conducted
in Denmark 10—12 years ago. His new work has the same merits that already distinguished
his older works: thoroughness, clarity and elegance of
presentation. Of the two volumes published so far, the first deals with
``det Gustavianska tidehvarfvets män'': Leopold, Rosenstein, Thorild
and Ehrenstein, whose philosophical ideas developed under the influence
of Locke and the French philosophers of the 18th century. The author seeks
to show that already these men bear witness that independent philosophical
thoughts could be born in his fatherland, and to refute the opinion that all
philosophical inquiry of any value in Sweden has been mainly an
echo from foreign countries. The second volume depicts Boëthius and
Höijer and their relation to Kant and Fichte. What remains is the depiction
of Atterbom, Biberg, Grubbe, Geijer and finally the conclusion of the whole movement
in the Boströmian system.

The philosophical movement proceeding from Locke had a decidedly
naturalistic character. It sought to carry through the method of experience
in all domains. By its clarity and boldness it became a mighty
weapon in the hands of the reformatory writers of the 18th century,
and just as its originator Locke was tied to the English revolution
of 1688, so it paved the way in the hands of those writers
for the French revolution of 1789. At the same time as the political
reaction, the idealistic speculation proceeding from Germany gained dominion
within philosophy, at least on the Continent. It is this
idealistic speculation that in Sweden took shape in a peculiar
manner. The place of the Boströmian philosophy in the history of philosophy is between
the schools proceeding from Kant and Fichte.

% --- p. 35 ---
\setcounter{footnote}{0}%
\opage{35}Although Prof. Nyblæus's work has not yet advanced to Boström,
we shall nevertheless, in discussing the two volumes at hand, be unable to
avoid discussing the Boströmian philosophy, since the author evaluates
the earlier philosophical movements by the principles of this philosophy, sees
it as the goal that is to be reached, and therefore regards the history of Lockeanism
and Kantianism as the self-dissolution of naturalism, its
successive transition to idealism. For one who has another conception
of the mutual relation of idealism and naturalism, the
history of modern philosophy will look different. It will therefore have
its interest to follow the author in his historical presentation; we
shall thereby come to discuss some of the most important principles in
philosophy. True, the author does not in the volumes at hand give the premises
for his critical evaluations; they will first be available in the last volume of
his work: the presentation of the Boströmian system. But we can
here provisionally orient ourselves with the help of Boström's own writings: besides
the one named in the heading especially his ``Grundlinier till philosophiska
statslärans propädeutik''.

A further contribution to enlightenment on how the history of philosophy
looks from the standpoint of Swedish speculation we receive in the
work by Boström's nearest predecessor, \textit{Samuel Grubbe,} which Prof.
Nyblæus has recently published. Its content is the history of ethics, but
it actually gives at the same time the whole history of philosophy. The presentation
is clear but somewhat broad. The author seeks to show the connection
and interaction between the ethical-philosophical domain
and the other domains of the spiritual and social life. Since the work
dates from the thirties and thus precedes the epoch-making works of Zeller, Erdmann
and others, it must be used with caution in
the details; but it bears witness both to reverence on the part of the editor and to
philosophical interest among his compatriots that an undertaking such as
the publication of the writings of this long-deceased investigator has been able to be
set in motion.

\begin{center}I.\end{center}

The popular consciousness rests securely in the belief in the absolute
reality of external things. It does not occur to it to make any distinction
between the objects of sensuous observation and things ``in
themselves,'' as they may be, apart from human apprehension, --- or
to trouble itself with asking how things ``in themselves'' must be
thought, or whether they can be thought at all. Just as little is
there for it any difficulty in assuming that our cognitions
arise through ``things'' acting upon us. When we learn something, we thus behave
as vessels into which the content of reality streams in.

% --- p. 36 ---
\setcounter{footnote}{0}%
\opage{36}It has cost philosophical reflection hard struggles to work its way
out of this way of representing things, which we are all compelled to
employ in practice, just as in everyday life we let the sun go round
the earth, however much we may be Copernicans in theory. And one
cannot wonder that it still asserts itself in many ways,
even among the thinkers who have studied most zealously the development
and essence of knowledge. It was John Locke who first made human
knowledge in earnest an object of study and thereby posed
the principal problems of modern philosophy. He did indeed show that the development
of knowledge rests on this, that consciousness, according to peculiar laws,
works over the particular representations drawn from experience; but in
the very reception of these single representations consciousness,
in his opinion, remained wholly passive. Prof. Nyblæus thinks that one can
describe Locke's theory as an effort to justify the
popular way of representing things, ``according to which man's knowledge is regarded
as a relation to an external object, an influence on our
consciousness by a reality lying outside it''. But what
right have we to ascribe such an effort to Locke? That he
at the beginning of his inquiries shares the popular way of representing things
cannot rightly be held against him, and one would on the contrary
with greater right be able to ascribe to him an effort to work his way
out of it. His distinction between primary and secondary qualities
(in which he has Hobbes, Galileo and Descartes as predecessors),
his half-skeptical conception of the concept of substance and his emphasis
on the activity of consciousness in the formation of the composite representations,
--- all this points decidedly beyond vulgar realism.
When one therefore speaks of the consequences of Locke, one has no right
without more ado to lay the materialism of the eighteenth century to
his charge and to say that only in ``Système de la nature'' ``the adherents
of Lockeanism learned to understand themselves and the content of the principles
which they hailed'' (Nyblæus, I, 1 p. 60). With the same right one could
say of Luther that he first learned to understand himself in Thomas
Münzer and John of Leyden. The author expresses himself still more strongly
in a following passage: ``They [Locke's French successors] brought to light
what the Lockean philosophy concealed in its womb, and showed that by
following its principles one comes into conflict with the whole higher consciousness of man
and with the very basic conditions of his life as a human being''
(p. 77). The question is precisely what Locke's principles were. He
is a Janus pointing in two directions, and the author himself says that it was only
by an uncritical and dogmatic procedure that the French materialists
arrived at their results.

The ground of Locke's great influence in the eighteenth century must
surely be sought less in the definite results of his philosophy than in the
analytic method, which he first applied with energy and on a larger
% --- p. 37 ---
\setcounter{footnote}{0}%
\opage{37}scale. For the ever stronger opposition this method was, as
Bouillier (``Histoire de la philosophie Cartésienne'' II, p. 550)
has remarked, ``an engine of war against the old social structure''.
As Locke had asked about the origin and elements of our knowledge,
so men now asked about the justification of the ruling social powers,
saw them as fruits of definite conditions and perhaps even sought to
imitate them. One cannot explain French materialism
from Locke alone, but must take the spiritual and social conditions of the time
into consideration --- an example of what we have emphasized in
the introduction: that one cannot always derive one philosophical school
from another, but must take account of factors that lie outside
philosophy. Materialism itself was a weapon that was seized upon as the
most defensible and most readily accessible. When one reads ``Système de
la nature'', one gets the impression of a religious and political pamphlet
rather than of a philosophical work; it has a decidedly polemical character. It
seems to me therefore entirely justified when Alb. Lange in
his ``Geschichte des Materialismus'' rejects the historical conception
current especially through the influence of the Hegelian school, according to which
materialism is derived from the philosophy of experience. It appears, after all,
even before Locke, in Hobbes, just as in France it appears
before Condillac, in La Mettrie.

Moreover Prof. Nyblæus himself intimates that the right development of
Locke's principles leads through Berkeley and Hume to Kant. What
mattered was to have the very concept of ``experience'' examined more closely.
But according to Prof. Nyblæus the method of experience is necessarily burdened
with the false presuppositions of the popular way of representing things,
and he therefore wants this method abolished altogether. History has,
in his opinion, also passed judgment on Locke and his method.
He describes the course of the history of philosophy in the eighteenth century in
the following way.

Berkeley and Hume showed that we have no right to posit
``things'' outside consciousness, since we have all along to do only with our
own sensations and representations. Knowledge cannot
without self-contradiction be conceived as a relation to external objects, since
these first come to be for us precisely in and through knowledge.

But besides this negative criticism an objection was also lodged
against Locke's empiricism from a standpoint that was indeed akin
to it, but in its ethical tendency rose above it.
There was a series of inquirers who indeed were not able scientifically to
master the principles of a higher and idealistic view, but
who yet had intimations of one. The author designates the school of these inquirers
as ``higher empiricism'', in contrast to Locke's empiricism
as the lower. ``Higher empirical realism has this in common with
the lower, that it places the Absolute in nature or in the reality given in space
% --- p. 38 ---
\setcounter{footnote}{0}%
\opage{38}and time. But whereas the lower empirical realism
reflects almost exclusively on the physical reality itself
or on matter (the stuff of which sensuous things consist), the higher empirical realism
on the other hand reflects on the \textit{formal} side of things,
on their \textit{determinateness of form,} or in other words: on \textit{order, coherence,
conformity to law} in the world and the principles of these'' (I, 1 p. 289).
The defect of the higher empiricism, too, is that it does not seek
the ground and final end of the sensuous (relative) world in an independent
non-sensuous reality (I, 1 p. 91). To this school the author refers
Rousseau, Herder, Lessing, and in Sweden Thorild and Ehrensvärd. But before the author goes on to describe the activity of these two men as thinkers,
he dwells first on the ablest representatives of Lockeanism in
Sweden: Leopold (b. 1756, d. 1829) and Rosenstein (b. 1752, d. 1824).

Whereas the latter relies entirely on Locke, Condillac
and Helvetius, \textit{Leopold} is properly a Lockean only in the theoretical respect.
From experience derives all content of knowledge, and reason
cannot go beyond it. But we possess another, higher sense
besides the one that is directed to the external: as the external sense
informs us of physical existence, so the internal informs us
of moral excellence. The moral-aesthetic feeling is a supplement
to theoretical metaphysics, and only thereby does philosophy
proper become possible, which is not merely a handmaid of the desire for knowledge, but
a protectress of the highest moral interests of mankind. Our
highest views cannot be made dependent on the results of the scientific
inquirer. It is contrary to duty to doubt the reality of the
truths inseparably bound up with the moral consciousness (God,
providence and immortality). To the ``natur-absoluta'' (i.e.
% sv: de yttersta vilkoren för all möjlig existens
``the ultimate conditions of all possible existence'') must therefore be added the ``moral-absoluta''
% sv: den högsta lagen för all menniskans fria verksamhet
(i.e. ``the highest law of all free activity of man''). (I, I p.
134 ff, 142—149).

In Leopold's character as a thinker the author finds an expression of the
Swedish national spirit, which also within the domain of philosophical speculation
is characterized by a predominant direction toward the practical
and a deep consciousness, connected with it, of the reality of personality
and of the supersensuous.

The harmony between theory and practice, nature and spirit, which Leopold
could not attain, was on the other hand for \textit{Thomas Thorild} (b. 1759, d.
1808) the content of the view in which his revolutionary spirit found
rest. Inspired by Rousseau and his gospel of nature he appeared
first as the champion of the new poetry (Young's, Klopstock's and Goethe's)
against the French school represented by Kellgren. Dissatisfied
that the government of Gustav III would not accept his proposal for
extended freedom of the press, he went to England to found
a world republic that was to be governed by geniuses and heroes. He accomplished
% --- p. 39 ---
\setcounter{footnote}{0}%
\opage{39}nothing for his purpose; on the other hand he made in London the acquaintance of
Swedenborg's writings. After his return to his native country he continued
his aesthetic and political polemic, and when after the death of Gustav III
in a writing addressed to the regent he gave vent to his enthusiasm for the
French Revolution, he was sentenced to banishment. Yet he received
an annual allowance from the government and a couple of years later became
professor at Greifswald, where he lived his last years and wrote
his most important philosophical writings (partly in German).

Thorild himself called his philosophy (in a letter) naturalism. He
wished to be able to say, when he died:
% sv: Jag såg naturen!
``I saw nature!'' But
``nature'' is then not taken here in the sense in which we took
this word in the introduction when we spoke of naturalism. By
``nature'' we understood the real given in experience (inner and outer), which
is subject to definite laws and develops under definite conditions.
Thorild, on the other hand, views nature as much poetically as philosophically;
his speculation is, as Atterbom\footnote[1]{``Svenska siare och skalder''. (Upsala 1847) IV. p. 12.} has put it: ``a poetry
affirming itself in thought''. He views nature as a great and
deep harmony, as a whole revealed in infinite, individual forms.
Reality is love in different degrees. All separate
beings are only different degrees or forms, in and through which
all-love develops its content, --- single tones in the great melody of the world
(``nur wegen der melodie sterben die tone'', he says
somewhere). Every single being can be understood only with a view to the
other beings, from which it differs only in degree; individuality
is ``quantum des \textit{quale}'', and it is the task of speculation to determine
``quantum dati'' in each single case, and thereby the place of each
single being in the great whole. Philosophy is ``archimetri'', i.e.
the doctrine of the eternal measure of the world; its task is \textit{``vide quantum!''}
% sv: känn förhållandet!
(``know the relation!'').

Abstract thought cannot apprehend the living, articulated
system of existence. Therefore Thorild polemicizes against Kant, Fichte and
Schelling, whom he regards as representatives of the abstract
school of thought which, instead of having its gaze turned toward the whole,
sees as little as possible and always the one without the other. There shall be
no divorce between thought and sensation; in true
knowledge the height of thought and the depth of sensation are united. Thorild demands
a kind of intellectual intuition, an ``allblick'', which by a ``götterblick
der allvergleichung'' embraces at once the most universal
and the most individual, or rather everything individual. He, he says,
who in his theory of the universe does not proceed from feeling, misses
the force, the soul, the Godhead in nature and sees in it only a dead
mechanism. Feeling, in which the abstract and the concrete have not yet
% --- p. 40 ---
\setcounter{footnote}{0}%
\opage{40}separated from each other, is therefore the first and last with Thorild.
Prof. Nyblæus characterizes his standpoint clearly by the remark
that whereas feeling in Rousseau had only psychological significance,
it acquired in Thorild ontological significance. It is therefore also in the highest
degree characteristic when Geijer says that the desiccated
materialism of the eighteenth century, which among the French and especially in Holbach wanted to make
lifelessness the general principle of existence, in Thorild turned into
a transitional system, which through feeling made everything alive and in this
all-vivifying feeling ecstatically practiced a kind of nature worship. It is undeniable
that feeling lent itself well to such a transitional system.
As feeling, namely, even when by its essential content it is purely
spiritual or non-sensuous, is always united with a sensuous element, so
we can understand how it could present itself to Thorild as the
reality sought by him, in which force and matter, the spiritual and
the bodily were united and fused''. (I, 1 p. 237).

When Thorild polemicizes against idealistic speculation, it is
thus not in order to assert the method of experience; he merely puts the authority of feeling
in the place of that of thought. But nevertheless Prof. Nyblæus finds too much
empiricism in him, especially in his demand that thought shall be controlled
by immediate observation and does not have its certainty in itself.
The astronomer's \textit{thoughts} about the starry heavens --- says Prof. Nyblæus --- are
after all something quite different from his \textit{observations} of them, and the latter could
not constitute the standard for the judgment of the former, as
Thorild from his standpoint would have to demand when he maintains that
``pure thoughts'' are only ``pure poems''. (I, 1 p. 217). --- But is it not
precisely a fundamental rule of inductive logic that thinking, however far
it may in its inferences and calculations remove itself from the immediately
given, must nevertheless always be held together with and compared with it,
and that knowledge of the real attains its completion only in the happy
meeting of inference and observation? That this completion
has been reached only at a few points is true; but astronomy precisely is a
magnificent example of how research, by relying on exact
observation, can arrive at a conception that turns the
immediate intuition upside down and yet agrees point for point with
it and explains it. Our knowledge is as a whole and in detail
still predominantly hypothetical, and the ground of this is precisely the difficulty
of verification at so many points. The sciences in which
subjective ideas can freely unfold, without being at once halted by
the strict police of verification, are farthest from their goal, and there is
no doubt that philosophy belongs among them. --- But the question of
the relation between experience and speculation we shall in a later context
return to.

That of Thorild's philosophical ideas which is of greatest significance for
later Swedish speculation is the thought of existence as an
% --- p. 41 ---
\setcounter{footnote}{0}%
\opage{41}organic system, a concrete and harmonious whole. This thought may
be said to be the fundamental thought of Swedish philosophy. But it appears
in Thorild not yet in the strictly idealistic form to which it later
worked its way. ``Had Thorild,'' says Prof. Nyblæus
(I, 1 p. 271), ``been able to give a consistent scientific development
of this idea of the whole, he would have been led to a rational-idealistic
world-view, or he would have seen himself compelled to
recognize that sensuous reality does not answer to the demands set up
for a whole, but presupposes another reality, which
must be purely ideal or purely spiritual, and therefore also purely rational
and personal.''

To a pantheistic world-view similar to Thorild's
was led \textit{Karl August} Ehrensvärd (b. 1745, d. 1800), whose significance
lies especially in the domain of aesthetics. ---

Whereas ``the higher empiricism'' could point to the ideal world-view only through intimation or poetic speculation,
a great positive progress was made with Kant's investigations of
the human faculty of cognition, which led to the result that consciousness
in its essence is active, unity in multiplicity, and that our knowledge
therefore indeed begins with sensuous experience, but yet does not
spring from it. Thereby an end was put to the presupposition
from which Locke proceeded: that consciousness behaved wholly passively in
the reception of the material of knowledge.

But Kant himself is, according to Prof. Nyblæus, still in part entangled in
the popular realism. He derives our sensations from the action of external
objects, and assumes a ``ding-an-sich'' outside consciousness.
Only the forms of knowledge have their original ground in the essence of consciousness;
the matter is taken in from without. Therefore he distinguishes sharply between
sensuous intuition and understanding: through the former alone objects are given, through
the latter they are only thought. Here, then, the presupposition seems to lie at the
ground that all reality must be given for sensuous intuition.

Kant has indeed, according to Prof. Nyblæus, acquired great merit by
sharply separating reason from sense perception and from the understanding that refers
itself to sense perception, and by conceiving reason as
the consciousness of the supersensuous he has ``uttered one of the thoughts by which
he leaves the empirical speculation of the preceding period far behind
him and opens a prospect on a new world''. But he denies that reason
by its own paths can attain any new knowledge, and holds that it merely sets
up an ideal, which demands progress beyond what is at any time given and
attained. Just as Kant, according to Prof. Nyblæus, is wrong in distinguishing
between matter and form, as far as the faculty of intuition is concerned, since everything that is apprehended
in the forms of intuition, as well as these forms themselves, lies
preformed in consciousness and has only to pass from possibility to
actuality, from obscurity to clarity, --- so he is also wrong
% --- p. 42 ---
\setcounter{footnote}{0}%
\opage{42}in conceiving reason as a merely formal faculty for setting up postulates
and ideals, instead of regarding it as a real faculty, which has the supersensuous
for its original content. Kant reverses the right relation,
``when, instead of seeking the cause of the need of human reason
to go back to the ultimate grounds in man's immediate
consciousness of the ideas, he on the contrary wants to ground the ideas on the
constitution of the human faculty of cognition and make them dependent
on it. This hysteron-proteron too contributes to Kant's rejection
of the constitutive validity of the ideas, a rejection that has its ultimate
ground in certain empirical and realistic presuppositions''. (I. 2 p. 35---37).

As in his doctrine of reason, so Kant has also in his ethics,
his philosophy of religion, his ``Kritik der Urtheilskraft'' given many intimations
of a higher world-view than that which follows from his theory of knowledge.
But whereas the side of his theory turned against realism is fully worked out, these intimations appear only
sporadically and receive no consistent carrying through in principle. Kant shows himself in
the determination of the most important concepts always infected with empiricism.
Thus e.g. in his consideration of the concept of the organism, which is so
significant for the ideal world-view. He proceeds here from
the empirically given reality, the sensuous organism, which develops
in time, and determines the general concept of the organism by it.
But the determination of the characteristic mark of the organism as self-development
fits only the finite or relative, not the perfect or
absolute organism. If, on the other hand, the unity of the organism is thought of as having
perfect power over its whole content, then there can be no question of
any development, since all development consists in making the
potential determinateness of the determinations through the unity actual. Development and motion
therefore do not lie in the concept of the organism when this is thought of without
limitation or in its truth and perfection. Thought of in
this way the organism means nothing other than a real system or
a reality that forms a whole closed in itself and fully coherent,
i.e. a reality whose unity is present in a multiplicity
of determinations or moments, which conversely are present in
the unity and in one another. (I, 2 p. 120—121).

There must therefore, according to Prof. Nyblæus, be distinguished between a lower
and a higher side in Kant. According to the lower side he shares Locke's and
empiricism's presupposition of an external reality, by which knowledge
is to be guided, and by which it is affected. The higher side consists
in the recognition that all reality is not exhausted in the sensuous world,
but that there is another and more perfect reality, raised above space and time,
than the sensuous one, which is only a phenomenon of it.
--- It is only the lower side that has been developed by Kant's successors
in Germany. Fichte indeed abandoned the assumption of a ``thing-in-itself''
outside consciousness, but he conceived nevertheless, like Kant (and
% --- p. 43 ---
\setcounter{footnote}{0}%
\opage{43}empiricism), consciousness as an empty form, which must first produce
its content. But production is motion, development, and presupposes
therefore time; consequently all spiritual life is conceived as finite.
Although Schelling and Hegel sought to remedy Fichte's one-sidedness,
their speculations nevertheless suffer from the same essential defect as his;
neither of them can conceive spiritual life otherwise than as developing
itself successively, and therefore finitely and relatively. ``Certainly it is
a fact that human consciousness develops its content successively
(in time); but this expresses only a defect or limitation
in us and consequently cannot lie in the nature of consciousness itself as
such or have any validity for the consciousness corresponding to its concept
(the absolute personality).'' This was first perceived by Swedish inquirers
(Grubbe and Boström). (1, 2 p. 122—127).

Before we follow the author further in his account of how Kant's
philosophy was received and worked over in Sweden, we must interpose some
remarks occasioned by his criticism of Kant.

There is one thing in which Kant is entirely at one with empiricism, namely
this, that all our knowledge must be drawn from experience. When thought wants to
go on its own, it behaves like the dove, which, because its flight goes
briskly in the atmospheric air, believed that it would go even better in
empty space. But it is the very concept of ``experience'' that Kant subjects
to a test, which leads to the result that all our experiences
are in their form determined by the nature of our consciousness. By analysis
of our experiential knowledge he sought the laws of the functions of consciousness,
--- which laws by the nature of the case must hold for all the content
that is developed and worked over through these functions. Whereas
Locke had acknowledged the activity of consciousness only in the formation of
the composite representations, Kant demonstrated such a psychological activity
in the formation of the most elementary sense images.

That the quality of our sensations too, not merely the forms
in which they are ordered, rests on our organization, Kant did not need specially
to impress, since it had been brought to provisional clarity by the earlier
philosophers' doctrine of the secondary qualities. When he nevertheless distinguished
between matter and form, so that the matter was given, whereas the form
is due to our activity, it was surely not his meaning that
anywhere in the course of the development of consciousness a wholly unformed
matter could be demonstrated. When consciousness --- as he taught --- is in its
essence synthesis, a combining into unity, then this characteristic mark
must hold for everything that emerges in consciousness. Kant's theory
has here been confirmed by the recent physiology of the senses. From the very
hour of its birth every sensation is brought into connection with other
sensations and determined by them. Matter in and for itself then becomes
merely something to which we approximate, the simpler the composition of consciousness
becomes. Pure matter is (as one of Kant's most acute
% --- p. 44 ---
\setcounter{footnote}{0}%
\opage{44}disciples, Sal. Maimon, expressed it) a differential, which is not met with and
cannot be met with in any real consciousness. Every real sensation
always has a certain form. In order that there may be consciousness,
several sensations must be given; but they must not merely be
as several; only through their unity in memory does consciousness exist.
Every moment of consciousness always combines several moments.

When Kant nevertheless says that an object becomes given to us only
by affecting our faculty of representation in a certain way,
he lets himself be misled by the popular mode of expression into confounding a
fact with its explanation. The fact is that we see, hear, etc.; but
that what we see and hear should also exist, even if we or others do not
see and hear it, of that we can of course know nothing whatever. But
Kant could also without harm to his theory have kept to
the simple fact that something is \textit{given} to us, that we see, hear, in general observe
it. This is the fact which the theory of knowledge has to start from and to
analyze. The real is what we apprehend as real, --- what we, however
much we may shrink from it, must in the end let stand as
it is, --- what we cannot but acknowledge. But the fact that we
``cannot but'' is evidently a purely subjective criterion, and no
other can by the nature of the case be given. For the dreamer
his dream is reality; but when he wakes, he discovers that his
dream was only an illusory reality, which was determined by a more
comprehensive reality and finds its explanation in this. No single
sense observation and no particular act of thought teaches us to know reality;
for that is required a constantly renewed and never concluded combining
of old and new experiences. As far as we can go in dreams
without running into jarring contradictions or conflicting experiences, so
far we also believe in the reality of the dream. But there comes a
point where the thread snaps. Even the most systematic dream is only
a fragment in relation to the totality that our waking experience presents
to us. And in a similar way it goes with those of our ideas that
do not have roots in reality; sooner or later their limitation will
show itself, and it is discovered that there is more in heaven and on
earth than we dreamed of in our philosophy.

The development of our knowledge is thus a progressive transition
from dream to reality, and our picture of reality is always only
relative. The more new things become ``given'' to us, and the better we understand
how to penetrate and analyze the given, the better we know reality.
But is there then no absolute reality from which all
that is given to us derives? Here we come over to a new
domain, where metaphysical theory takes the place of logical-psychological
analysis; and in order to have firm ground under our feet it is important
not to confound these two things, as the popular
conception and many philosophers do.

% --- p. 45 ---
\setcounter{footnote}{0}%
\opage{45}It will thus be seen that one can cut away the popularly realistic presuppositions
from Kant's theory without thereby immediately hurling oneself into
the arms of idealist speculation.

There was also another point in Kant's philosophy where he sinned
against his own principles. Just as, by directly positing a thing in itself, he went beyond experience and consciousness, so
he let, at one point, namely in the command of duty, absolute reality
step forth directly into consciousness. But since the command of duty is a psychological
fact, it must hold of it, as of all other
facts, that it is a phenomenon subject to the laws of experience
(here of psychological and historical experience). Apprehension
of an absolute reality is, according to Kant's own philosophy, a
\textit{contradictio in adjecto,} since reality, by being apprehended by us, is \textit{eo ipso}
relative reality. Here too, then, Kant goes behind consciousness,
--- and the example he had given of breaking the principles of the critical philosophy
was faithfully followed by his successors.

It seems to me that Prof. Nyblæus, by his distinction between
a ``lower'' and a ``higher'' side in Kant, has come to overlook
what is central in him: the critical standpoint, in which after all his
true greatness and significance lie. To give profound allusions
to an ideal world, mystical pointers to a ``higher reality,''
is, after all, not so difficult; one has the whole Platonic and Neo-Platonic
literature at one's disposal, not to speak of theology and
the catechism. But to become the true founder of the theory of knowledge,
to lay bare by logical-psychological analysis the elementary
structure of the human spirit, that takes a genius and demands a labor of thought which
seeks its equal.

It is not my opinion, with certain zealous Neo-Kantians, to conceive
post-Kantian speculation as a fall from grace. In any case
Kant himself had put the serpent into paradise. In order to
be just to the heroes of speculation one must remember what abuses could
be made of Kant's philosophy, quite apart from its inner defects. There
is no question that what won so many for Kant's system
and so quickly gained it entry was quite as much his practical as
his theoretical results. It was the deep ethical seriousness that
captivated so many. But among them there were again not a few who consoled
themselves for the impossibility of any speculative theology, for the overthrow
of the good old ``proofs of the existence of God,'' with the direction
Kant himself gave to postulate God and immortality on moral
grounds. They found a new way out to save their old dogmas.
They set Kant highest precisely there where he ceased to be a
thinker; for the God who in Kant is to bring virtue and felicity into
harmony is just as mythical a being as the ``Demiurge'' who in Plato
is conjured up to form matter after the ideas. With justice
% --- p. 46 ---
\setcounter{footnote}{0}%
\opage{46}the young Schelling called out to such Kantians (in his interesting ``Briefe
über Dogmatismus und Kriticismus''): ``For you it was not dogmatism
but at the most only the dogmatic philosophy that was overthrown!''
--- It was thus in order to assert the freedom and independence of philosophy
and to satisfy the demand for unity in the world-view that so many
gifted spirits took to the speculative paths. When one would not
go by the path of experience, which does not lead beyond relativity, there stood
only one path open: idealizing abstraction.

When Prof. Nyblæus in his critique reproaches Kant for regarding
temporal development as an essential element in the concept of consciousness as well as of
the organism, although time holds only for the finite, not
for the infinite and absolute, he owes it to Kant first to give an account
of how it becomes possible to make such a transition from the
finite to the infinite and in general to set up a positive concept
of the absolute. To this critique Kant would quite simply reply
that he has spoken only of the organism and the consciousness that we
\textit{know;} only when a consciousness without conditions of time can be demonstrated, an organism
not subject to development, will it be justified to
extend the concepts. The question is whether we can at all think
the timeless, --- whether change, transition, movement, development
are not fundamental determinations in the nature of our consciousness. And even if
we could abstract from these determinations, the mere abstraction
is still not enough to form the representation of a \textit{real} being.
By the path of abstraction we can get back to the simple fundamental relations
and fundamental laws, but it cannot lead to any new knowledge of
a real that is not already given. To appeal to ``man's
immediate consciousness of the ideas'' (I, 2 p. 36) helps nothing, since it
is surely not everyone who possesses such a consciousness, and it must in any case
be submitted to a closer psychological analysis before one can
know what it means. If human reason possessed such
an original content and were a ``real faculty,'' we should get two kinds of
logic: one for the higher and another for the lower experience. --- But the
ultimate motive for setting up the ideal of the perfected reality, raised above time and
development, is --- as will later appear --- in
truth not logical, but ethical.

It was also the path of abstraction that was taken in Germany when,
after the critical interlude, the speculative paths were trodden anew.
Kant's analysis had led to the conception of consciousness as an original
synthesis. But synthesis always presupposes a plurality that
is united, a duality that is overcome. Let us then --- it was said ---
go back to the point where the duality breaks forth, where the manifold
is not yet posited, --- to pure activity! In all real
consciousness we find, to be sure, only the relative, determinate and
modified activity; but for that very reason speculation must also build on
% --- p. 47 ---
\setcounter{footnote}{0}%
\opage{47}a ground that lies beyond consciousness. All real consciousness ---
said Fichte --- rests on a duplicity, a distinction between subject and
object. The pure activity through which subject and object first come into being
has therefore itself no object, no resistance. In order that consciousness may
come into being and experience become possible, limit and restriction must be posited;
only relative reality is reality for us. But why then should
experience be possible? To this question Fichte gives only a teleological
answer: because the infinite, expansive movement does not have the
value that consciousness and knowledge have. Becoming conscious of oneself is thus a
good that is to be attained, and the eternal process of reason aims at that. In
different forms it is the same answer that Schelling and Hegel too give
in the last instance. Fichte is of such great interest because the motives
and starting points of speculation lie so plainly in him.

We stand here at the point where the ways part. The path that leads from
Kantianism into the land of speculation: idealizing abstraction,
divides into two branches, as German speculation boldly ascribes
to the absolute an eternal self-development, whereas Swedish speculation
regards the absolute as the eternally perfected and completed.
We will not yet inquire which of these schools has
the greater value and justification; for we still have to consider the
historical development of Swedish speculation.

The Kantian philosophy met in Sweden with sharp criticism from several
of the men who had joined the school of Locke; besides Leopold
there appeared especially \textit{Fremling} (d. 1820) with an exceedingly acute
critique, above all of Kant's moral theology. He shows how Kant here
constantly comes into contradiction with himself: The postulate of immortality
Kant grounds on the claim that man here in this life never attains
ethical perfection, --- but this, as he himself teaches, cannot
be attained by any finite being, but remains standing as an ideal. Kant
does say that time has no validity for the intelligible world;
but how does he know that, when this intelligible world cannot be
an object of knowledge? A personal immortality does not in any
case follow; for it might be that perfection were realized
through the development of the race, which is gathered by the Godhead into unity.
And when Kant in his Ethics would abolish self-love, the
wish for the continuation of personal existence seems also bound to be pushed
back. It is finally one of Kant's methodological principles
that everything in the universe is to be explained by natural causes, so that even the
wildest physical hypotheses are to be preferred to an appeal to the supernatural;
but then he seems also, in the question of the connection of felicity
with morality, bound to content himself with a moral order of the world
conditioned by natural causes.

Other adherents of Locke, by contrast, were won over to the new
philosophy. Kant's first disciple in Sweden was Daniel Boëthius (b. 1751,
% --- p. 48 ---
\setcounter{footnote}{0}%
\opage{48}d. 1810). It was through this man's teaching at the university
of Upsala that idealist philosophy gained entry into Sweden and
the impulse was given to the national philosophical school. Not only
Höijer, but also Biberg and Grubbe received from him the first awakening
to philosophical inquiry. He gave the first sketch of the
ethical-religious world-view that was later perfected by the two last-named
% sv: Utan den djupa sedliga anda, som utgick ifrån en Boëthius, en Biberg, en Grubbe och Boström, skulle säkerligen icke ens Geijer kunnat bilda en tillräcklig motvigt emot det upplösande moment, som låg i ”fosforisternas” naturalism.
inquirers and Boström.\footnote[1]{Of interest for the history of general literature is the following remark by Prof. Nyblæus (I, 2 p. 236): ``Without the deep moral spirit that proceeded from a Boëthius, a Biberg, a Grubbe and Boström, not even Geijer could certainly have formed a sufficient counterweight to the dissolving element that lay in the naturalism of the `Phosphorists.'\,''}

While Boëthius agrees with Kant that reason cannot theoretically
attain the absolute, he lays more weight than Kant on the fact that
it demands harmony in our essence. This thought of harmony, of the
spiritual life as an ideal organism, is --- as we have already
seen --- the alpha and omega of Swedish philosophy. By means of it Boëthius is led
beyond the Kantian opposition between theory and practice, between reason
and feeling. He seeks to give a genetic explanation of the categorical
imperative by seeing it as the form and expression of thought's self-activity,
from which all law and all harmony spring. He is thereby led to
a fusion of Kant's ethics and the Greek, in that he places the
ancient admiration for the morally beautiful (\textit{τὸ καλόν}, \textit{honestum}) alongside
Kant's doctrine of respect for the law of reason. At the same time he conceived the
ethical as in its ultimate ground one with the religious, morality as
``a life in God.''

The real chief representative of the Kant-Fichtean school
in Sweden, and on the whole one of the most important Swedish philosophers, is \textit{Benjamin}
Höijer (b. 1767, d. 1812). He was, like Thorild, an ardent
friend of freedom, and when he had joined a literary club in Upsala
called ``Juntan,'' which was accused of Jacobinism, he suffered
on account of
% sv: skuggrädslan hos Gustaf IV Adolf
``the shadow-fear of Gustaf IV Adolf'' a series of humiliations
and slights and was appointed professor only after this
king's deposition, but soon afterward he fell into an illness, the effects of which
he never overcame.

``Höijer was,'' says Prof. Nyblæus (I, 2 p. 351), ``one of those superior
natures who feel called upon to give laws in the world of knowledge
to the rest of mankind. --- --- With the superior intelligence
and the strong self-esteem (which, however, overstepped the bounds of
propriety) there were united in him an undaunted courage, which did not hesitate
to draw the consequences of the principles that he considered
right, and a firm will, which did not tire in the endeavor to reach
the goal he had set himself. Characteristic of him in the first respect
is an answer that he is said to have given to a well-meaning person,
% --- p. 49 ---
\setcounter{footnote}{0}%
\opage{49}who thought he ought to draw his attention to the dangers to the
salvation of his soul with which free inquiry was bound up:
% sv: Sök sanningen, och bure det än til helvetets portar, så klappa på!
``Seek the truth, and though it lead to the gates of hell, knock on!''

On a journey in Germany Höijer had made the acquaintance of Fichte and Schelling,
and both these German philosophers rated him very highly as a
thinker. His position in the history of philosophy is characterized by the fact that he successively
anticipates the ideas of Fichte, Schelling and Hegel and thus for his
own part describes the same course which in Germany, with
wider execution and richer fullness, was traversed by those three thinkers.
He saw the necessity of a more inward coherence and unity
in philosophy than Kant, whom he otherwise regarded as the greatest of
all philosophers, had achieved; and he wished, like Fichte, to conceive
consciousness not merely as the ground of the form of knowledge, but
also of its content. But Fichte remains standing at the abstract
self-consciousness, the pure I, and therefore does not get beyond dualism,
since an I is unthinkable except in opposition to an object. Therefore
Höijer goes back to the principle that is common to subjectivity
and objectivity, and thus anticipates Schelling's doctrine of identity.
Whereas for Fichte it was impossible to demonstrate how the pure I
comes to posit limits for itself and thereby become conscious, Höijer
seeks to demonstrate limitation as originally given in the very concept
of action. Without concentration no energy; without limitation
the absolute would lose itself in infinite expansion. But he did not
agree with Schelling in the position the latter assigned to the philosophy of nature as
coordinate with transcendental philosophy; nature was for Höijer, as
for Hegel, only a sphere of transit for the absolute. Nor
could he rest content with Schelling's ``intellectual intuition'';
he demanded, like Hegel later, a grounding of this intuition
by a higher logic; and it was works in this direction that
occupied him in his last years.\footnote[1]{In his last years Höijer regarded with growing dissatisfaction the way in which the younger Swedish Schellingians conducted themselves in their journal ``Fosforos.'' ``Although Höijer's attitude toward them was originally merely evasive, but far from hostile, and more waiting than rejecting, yet the more they betrayed what they had in their shield, he could not avoid fearing that the kind of Schellingianism which they sought to spread in our country would work destructively upon all truly philosophical study at the university and bring philosophy into complete disrepute. It was probably chiefly from fear of this that Höijer wrote a treatise on `åskådning,' which is found in volume 2 of his writings, and in which he castigates the `philosophy of words and formulas into which Schellingianism had degenerated among the multitude of its adherents'. (I, 2 p. 349). --- Even on his deathbed he disclaimed with great zeal any part in the said journal.} Finally he could not, like Schelling,
regard art as the highest form of expression of the spiritual
life; for him, as for Kant and Fichte, practical-ethical
activity stood highest.

According to Höijer it is a practical interest that leads to philosophy,
the fundamental philosophical problem being how the moral
% --- p. 50 ---
\setcounter{footnote}{0}%
\opage{50}consciousness and the consciousness of freedom inseparably bound up with it can be preserved
in the face of natural necessity, or --- seen from another side ---
what the relation is between subject and object. Therefore reason
is not satisfied until it has reached the absolute, where subject and object
are one. The demonstration of the reality of our knowledge and the explanation of
the possibility of objects coincide in the philosophical construction,
which derives relative reality from the absolute and shows how
the finite I's, for whom the object or nature asserts itself,
are nothing other than the multiplicity of self-limitations
under which the absolute reveals itself, and by which it is completely
determined. --- Human knowledge must be a coherent
whole in itself, a system; and in order to be an object of knowledge
existence itself must also be a system. This is Höijer's main thought,
and by its further working out he arrived at the result: that
since the objects of the external world must be conceived as unconscious products of
the operation of the absolute, one may justly say that nature in its ultimate
ground is nothing other than the moral world seen from a certain side, an
expression of man's place within this world. If the
human individual did not belong to a whole of beings who recognize
morality and advance toward it, there would be no moral perfecting
for him; for only within such a whole is duty (i.e.
% sv: allmän lags giltighet
``the validity of a universal law'') possible. (I, 2 p. 358, 364 ff, 401 f).

Höijer still belongs, like the German thinkers most closely related to him,
according to Prof. Nyblæus, to those who have developed the
``lower side'' of Kantianism. He applies the concept of development to
the absolute, lets it ``come to itself'' by overcoming
limits and oppositions. In reality he has therefore not got beyond
empiricism and realism, has not reached the concept of
true reality! (I, 2 p. 443).

Unfortunately we lose at this point Prof. Nyblæus's expert guidance
and must on our own and at our own risk add some traits
to the characterization of Swedish idealism in its perfected form. The
most important principles we have already met sporadically in what precedes.

Speculative philosophy seeks the absolute and finds rest only in
it. The different speculative schools therefore part from one another
by their different conceptions of the absolute. ``The highest
concept in the whole of philosophy,'' says Prof. Nyblæus (I, 2 p. 424), ``is the concept
of the absolute or of independent and perfect reality,
to which the finite world, on account of the lack of independence
and the imperfection with which it is in a certain respect afflicted,
points as its ground of explanation.'' --- But Swedish idealism,
which, prepared by Biberg, Grubbe and Geijer, received its completed form
from Boström (b. 1797, d. 1866), finds decisive deficiencies in most
earlier attempts to develop this concept. Absolute, says Boström,
% --- p. 51 ---
\setcounter{footnote}{0}%
\opage{51}is that which is separated from and independent of everything else, and which in itself
contains all that constitutes its essence. Absolute is thus the same
as perfect; no determination can be taken from it or added
to it. Absolute beings therefore cannot have anything outside themselves,
cannot be subject to time and change, space and place. The form of
self-consciousness is the only one possible for an absolute being; for in perfected
self-consciousness there is one in all, and all in one (\textit{πάντα ἐν παντί}, as
Anaxagoras said). (``Om relig.'' etc., p. 88—94, 32---34). As an independent
and perfect totality the absolute possesses \textit{absolute} infinity;
its perfection is always real, and it is impossible to go
further. The sensuous, by contrast, possesses relative infinity, in that its
perfection never becomes real, but it is constantly possible to
go further. At the same time the absolute possesses \textit{inner} infinity, i.e.
none of its moments can be conceived by itself alone, none can be thought
without all the others being \textit{eo ipso} thought; they are reciprocally one another's
moments. The absolute reason has in its moments subjects and
centers for the apprehension of itself and senses itself with its whole
content in the complete sensation in any one of them whatever. Thus it has them all
present in itself, and it is itself with its whole content present in each one among them; whence it
follows that they are also all present in one another. (``Grundlinier''
etc., p. 30). All moments do not occupy the same rank; there are lower
and higher beings in the absolute. But the one is not outside the
other; the relation between them can be symbolized by the relation between the
numbers: 99 is less than 100, but 100 includes 99 in itself. The
higher beings include the lower and in addition a plus. Thus
the individual persons enter as moments into society as a
more comprehensive individual, and the smaller societies into the higher. (``Grundlinier,''
p. 40—41).

Here, then, we have the perfect organism, the ideal harmony
described, to which we were so often referred in what precedes by the historian of
Swedish philosophy. It is a \textit{mundus intelligibilis,} a
spiritual universe, that is here presented to us as the true reality.
It is an absolute personality --- but precisely as an absolute personality
it is a realm of personal spirits, a system of absolute beings.
It has its content originally in itself; the content cannot arise, cannot come
from outside, but must be as eternal as the absolute itself.

When the absolute is thought in this way, all those conceptions
must be rejected which let its content arise through creation or development
and thereby make it finite. When theology lets God create the world,
it makes him a finite being, since creation must after all produce
a change in his nature, and he thus could not beforehand be
absolutely perfect. When German speculation lets the content of the world
arise through a process which is an evolution of the absolute's own
% --- p. 52 ---
\setcounter{footnote}{0}%
\opage{52}essence, it likewise brings finitude into it and thereby entangles itself
in contradictions. Both the theory of creation and the theory of production are still infected
with empiricism, in that they regard movement as an essential element
in the nature of life and consciousness. It is precisely a sign of relative
deficiency or negation of life when living beings move. The being
that is in all respects living and self-conscious cannot and need not
develop and change, as the finite must do. It is
only independence in movement that essentially belongs to life; but
perfect independence is found only in the absolute. (``Grundlinier''
p. 1—2).\footnote[1]{The polemic against ecclesiastical theology and German speculation is found in the two sharp polemical pamphlets: ``Anmärkningar om Helfveteläran'' and ``Den speculativa filosofen J. J. Borelius i Calmar.''}

If one wished to call such a view pantheism, Boström replies
that by this word one usually understands a
doctrine which does not think of the Godhead as independent and autonomous
of the sensuous world, but represents the former as the latter's substance,
and the latter as the former's accident, so that both mutually presuppose each other.
``According to my view, by contrast, the sensuous world can certainly not
exist without the Godhead, but the latter presupposes and needs the former as
little as, e.g., the real sun needs the visible sun. For
me the spiritual is namely the only real and originally existing;
and since I conceive the Godhead as the absolute spirit, I must conceive
everything truly existing in it and not set it in immediate relation to
anything external. The sensuous world, on the other hand, I apprehend as constituting only
a phenomenon for us of the spiritual and supersensuous, and
standing in about the same relation to this as, e.g., twilight
to daylight, or the visible sun to the real. --- If this
can be called pantheism, then it is in any case not the usual kind.''
(``Satser om lag och lagstiftning''. Andra uppl. Upsala 1871. p. 54).

There is thus another universe than that ideal one, infinitely
harmonious and perfected. But if we ask where this sensuous universe,
which is in time and space, which changes, develops and dissolves, is
to be found, and how it arises, --- then Boström answers that finite
beings, precisely on account of their finitude, cannot apprehend
the absolute (the ideal world) with absolute clearness and distinctness,
for which reason it appears to them as something other than
it is in itself. It is only the infinite reason or God who
can apprehend himself and through that everything else clearly and truly. All
other beings must sense themselves and through that everything else more
or less dimly and unclearly and to that extent relatively and imperfectly.
This has its ground in the fact that every other being than the highest is finite
and therefore in its content imperfect in comparison with
the higher and the highest: it has a lower measure of capacity of apprehension,
% --- p. 53 ---
\setcounter{footnote}{0}%
\opage{53}on account of its place in the system of beings. Nothing of the eternal content
is therefore lost to the finite beings, but a shadow is cast
over it; it appears in an unclear, sensuous form. The sensuous is
only the rational in a dim, potential, undeveloped (unseparated) form. Each
single finite being finds itself at its own peculiar position and has
therefore its own individual universe, which is more or less in agreement
with those of the other beings, according as there is greater or
less kinship in their nature.\footnote[1]{In a similar way everyone has in the ethical respect his own law and cannot follow any other (``Satser'' etc., p. 17---21), since every law that he recognizes is recognized by him because it expresses his essence, whether this be sensuous or rational.} There are thus as many worlds
as there are finite beings, and each world has its determinate
degree of clarity. The progress of knowledge consists in the dimness
gradually yielding, and the eternal and infinite content also being apprehended
as eternal and infinite. The only development that takes place is
this transition of the finite beings from dimness to clarity.
This development takes place through a series of ever higher forms of life.
But even in the highest form of life sensuousness --- and thus
dimness --- has not vanished. It is indeed less than at the beginning
of the development, but if it were completely abolished the
finite beings would lose their independent subsistence. The development of every being
stops when the highest perfection possible for it has been
reached; and there remains always a difference between the system of the finite
beings in their highest forms of life and the system of the purely ideal beings
or the ideas, as they are thought by the Godhead itself and
are its self-conscious determinations. (``Grundlinier'' p. 44 ff, 64 ff).

One will easily see that, as regards the finite beings, Boström
acknowledges the law of relativity; for each of them there exists only
a relative reality, and they can in all eternity not get beyond
this relativity. On the other hand he assumes that the case is otherwise
with regard to the absolute being: it sees things as
they are in themselves. But does it not lie directly in the very concept of
knowledge that this is impossible? Knowledge presupposes a reciprocal determination
of subject and object; but determination means relation,
and that which stands in relation to another is different from it. This
Boström himself states. ``It is clear,'' he says (``Grundlinier'' p. 38),
``that nothing can be, and be something in and for another, unless it
first (in the concept) is, and is something in and for itself; for that which was nothing
in and for itself could not be thought to stand \textit{in relation to}
anything else --- which is necessary if it is to be thought as one of
its determinations.'' If, then, relation cannot be kept out,
can relativity be kept out? Since knowledge is always determined
by the nature of the knowing being, then God's knowledge too must be
determined by his nature; he must have his phenomenal world, just as
% --- p. 54 ---
\setcounter{footnote}{0}%
\opage{54}we have ours, and he can --- as Jouffroy has very rightly remarked\footnote[1]{``Dictionnaire des sciences philosophiques'' 2:e éd., p. 830.} ---
know just as little as we do whether his knowledge agrees with the
absolute reality, or rather: the ``absolute reality'' remains,
however high we may soar in the series of real or thought beings,
with however ideal attributes we may equip them, always only an
ideal, a limit that constantly recedes.

The fact is that, while Boström has clearly seen the relativity of concepts
such as time, space, force, development, etc., and therefore seeks to keep
them away from the absolute, he has not seen that the concepts of knowledge and consciousness
are likewise relative, and that they are therefore canceled
when one tries to give them absolute significance.

But a no less serious difficulty arises when for a moment we
acknowledge the reality of the ideal world, --- the same difficulty that
was already present for Plato, and which idealism has never been able to
overcome. Whence, namely, comes the double (properly infinitely
manifold) mode of apprehension of one and the same universe? Boström
answers that it is finitude that causes the dimmer apprehension.
And finitude itself? Boström answers (``Anmärkningar'' etc., p. 14):
``To ask for a cause of the finitude of the human spirit would be about
as unreasonable as to ask for a cause of the three-sidedness of the triangle.
It would not be a triangle if it were not three-sided;
and likewise the former would not be a human spirit if it were not
finite.'' Thus one may here no longer ask \textit{why;} we stand
at something that can no longer be grounded. The lower measure of capacity of apprehension
that the finite being has rests on its ``eternal place
in the system of beings.'' (``Grundlinier'' p. 45). Beyond the opposition
between the infinite and the finite we could thus never get:
it is eternal.

And yet it is in order to explain the sensuous world and finite
development that thought has constructed for itself that ideal world! It
is a maxim in Boström that the higher explains the lower (``Grundlinier''
p. 65 f), that the perfect explains the imperfect, so that there
could not be imperfect knowledge unless there were
perfect knowledge. (``Om relig.'' etc., p. 47). But have we in
reality gained any explanation by going back to the ideal world?
We have only doubled the problem, just as Plato thought to
explain real things by forming an idea for each of them.

Perhaps one will hold to Boström's image: that the sensuous world
presupposes the ideal world, just as the visible sun presupposes the
real sun, and twilight presupposes the full light. But this
image seems to me to be highly ominous for the Boströmian
view. For there would be no difference between the sun as
% --- p. 55 ---
\setcounter{footnote}{0}%
\opage{55}globe and the small disc we see in the sky, unless besides the solar globe
there were another globe from which it could appear as a
small disc; but where is such a position possible, when the ideal world is
the only reality, and there is nothing besides it? And in order that light
may be dimmed, a positive obstacle is presupposed; but where is there
such a positive obstacle for the absolute? Boström indeed lets
twilight be a unity of light and darkness; but darkness is nothing
other than precisely the negation of light --- so one is none the wiser for that. When
he says that ``the phenomenal world is a result in and for us of the cooperation of two opposite
factors, a positive and a negative'' (``Grundlinier'' p.
46), this too gives no explanation, since the absolutely negative
(= nothing) cannot act at all, and we lack information as to what
positive factor it is that in this particular connection can be designated
as negative, because it acts in opposition to another factor. But
such an independent positive factor would conflict with the presupposition
of the ideal world as the absolute reality.

Swedish philosophy leads us step by step from the low regions of
empiricism upward to the mountain peak of speculation. Even in Hegel, the
philosopher of panlogism, it still finds ``a good deal of empiricism''; he is said
not yet to have overcome the after-effects of Locke. It is on the
whole characteristic of the historical position of Swedish speculation that
it constantly conceives itself as the opposite of the course of development from
Locke to Hegel, --- as the Elysium where the thoughts that fought during this long period of
development may rest in undisturbed harmony. The question
is whether by this strict idealism it does not precisely work against
itself. It seeks absolute guarantees for human knowledge
and action, and no one can fail to recognize that a theoretical
and practical ``excelsior!'' resounds through it. But it ascends so high that
it loses sight of reality --- of the reality which was after all precisely
to be explained; it leads us ever higher up, but thereby arrives
at last in the regions of eternal snow, where earthly growths no
longer thrive. It is only in ``the fertile lowland of experience'' that we can
dwell and build; even if we now and then climb up the heights to orient
ourselves and gain an overview, we still do not build our houses up there.

\begin{center}II.\end{center}

After the historical survey now given, we will, for the further grounding
of the critical remarks put forward on various occasions, go somewhat
more closely into some particular questions that are of special importance for the
conception and treatment of the philosophical problems.

Even the first attempts to present the method of knowledge
acknowledge that this cannot be single, but that it consists in a
double movement. Thought finds its content and itself as
% --- p. 56 ---
\setcounter{footnote}{0}%
\opage{56}immediately given, and seeks now to order, group, dissolve and combine
the different parts of this content. It finds a chaos and would produce
a world. The first movement is an analysis, by which the
single elements are separated out from their first, accidental connections;
thereupon synthesis seeks to bring about new connections according to the elements'
kinship and essential belonging together.

Knowledge is a struggle with the unknown, and originally everything
is unknown. The unknown is overcome by being stripped, layer by layer,
of its immediate form and clothed in a dress with which thought is familiar.
There is in itself no fixed goal set for analysis; the possibility
of proceeding further is constantly present. One halts when the
given has been traced back to relations so simple and transparent that
for the time being it is not possible to get further. Synthesis then seeks to compose
the elements into a totality. Analysis goes from things to
grounds, synthesis from grounds to things. In the interaction between
both consists the life of science. Goethe has aptly compared
them with inhaling and exhaling. The ideal of science is a complete
carrying through of analysis, whereby in turn a complete synthesis
becomes possible. The rich and manifold content of reality would
then be traced back to its simplest expression, to a formula in which its
whole fullness was contained.

This duality in the method will be found grounded, not merely
by an investigation of the special methods of the different sciences,
but also by a consideration of the general essence and
conditions of knowledge. In all life of consciousness there is an indissoluble union of unity
and multiplicity, of activity and passivity. The scattered, isolated
states become conscious only by being brought into connection with one another:
no sensation is absolutely simple. There are no absolute
atoms of consciousness. But unity is at the same time an activity, which makes
itself manifest in the incessant --- voluntary or involuntary --- comparison,
distinguishing, juxtaposing, dissolution and connection, which goes on
in consciousness. The single impression that we receive is drawn, so to
speak, into a maelstrom immediately upon its entry into consciousness. The
place that is assigned to it in the circle of representations depends not only on
itself, but predominantly on its relation to other given or reproduced
representations. One will now easily see that in analysis attention
is directed to the single determinations and seeks to apprehend them
as sharply as possible, each by itself; in synthesis the interest prevails
in restoring the unity again and with it the subject's own inner activity.
There come to be, then, two poles in the world of knowledge. And this agrees
furthermore with the significance of knowledge for the individual. From the first
it is no purely theoretical concern. It is a question of
life and death whether the individual orders and connects his representations in the
right way. However much knowledge therefore always retains its
% --- p. 57 ---
\setcounter{footnote}{0}%
\opage{57}subjective character, it must constantly strive to get beyond this,
strive to bring its connections of representations nearer to objective
validity. The life of knowledge is only a single form of the great
struggle for existence and of the demand that is made upon all organic
beings: that they bring themselves into harmony with their environment. An
isolated subjectivity would be the sure prey of destruction.

Analysis is everywhere the foundation of synthesis, but the closer relation
between the two methods --- whether a synthesis is possible at all, and
to what extent --- is different in the different sciences.
The synthetic method can acquire positive significance only when
analysis has led to simple and uniform relations that in themselves contain
the principle of an entirely new and independent development. This is
thus the case with the analysis that leads to the setting up of the mathematical
axioms. In other cases synthesis becomes only a reconstruction
from a certain point, with constant regard to experience and
with successive taking up of the elements given by analysis. Reality
is always coherence, concrete totality. But the immediate
coherence appears to thought as accidental; it then depends
on whether it can, as the fruit of its labor, set a new totality
in place of the given. It often fares with it as with children, who tear
their toys to pieces to find out how they look inside,
and now try in vain to put them together again. And yet
children are entirely right that they cannot get to know things
as long as they let them remain as they are.

Philosophy must be subject to the same law as all other branches of science.
There are no shortcuts to the end of knowledge; the way is
everywhere one and the same. Speculative philosophy has nevertheless often
taken its starting point in general principles and ideas whose validity
it did not doubt, and whose consequences it then pursued as far as
it could. But one will always find that an analytic
work has gone before --- as indeed Socrates goes before Plato, Descartes before
Spinoza, Kant before Fichte, Schelling and Hegel. The fault here is only
that one has too quickly considered the analysis complete. On the
other hand the speculative systems themselves must be regarded as different
attempts to investigate particular fundamental concepts and test their
range. The single concept that is made the principle of the system is
taken out of the whole human circle of ideas, and in this way
the history of the systems is a progressive analysis, by which the views
are clarified and purified. Yet as a rule the individual system shrinks from
thus being a link in a progressive process; it believes that in its principle
it has penetrated to the interior of existence and demands that it be recognized
as the true knowledge. Thereby philosophy so often acquires a dogmatic
character, and it will always be found that dogmatism is connected
with an underestimation of the analytic method. As synthetic
% --- p. 58 ---
\setcounter{footnote}{0}%
\opage{58}science philosophy is system, as analytic science critique. On
the relation between the systematic and the critical depends a philosophical
attempt's character and value. But once consciousness has awakened
to the significance of the principles themselves being investigated and grounded,
one will soon see that the scientific character of philosophy rests on the
critical-analytic method. The old rule de principiis non est
disputandum is, when it is to hold absolutely, wholly reprehensible. That
one must always in one's thinking proceed from certain principles is certain;
but these principles themselves it may in many cases be
fully justified and necessary to draw into the discussion.

When we wish to examine more closely how philosophy stands
with regard to the relation between analytic and synthetic method,
we must distinguish between the different disciplines that are commonly
reckoned under philosophy. Whereas in antiquity all knowledge was called
philosophy, there have since gradually --- as by a cell formation out of
the structureless protoplasm --- one special science after another formed
and secreted itself. In English physics is still called ``natural
philosophy,'' in memory of its tender infancy. The most recent
example of such a secretion is political economy, founded by Adam Smith.
Some of the disciplines that are still reckoned to philosophy
have, like the special sciences, an entirely special and limited
domain, and they will therefore in time be able to constitute themselves as independent
sciences. Thus psychology, logic and ethics; they seem
almost to stand on the verge of it. When in this way all particular domains each get
their special branch of science, their analysis and synthesis, there remains only
one task, which falls to philosophy in the general sense of the word
—: that of investigating what value and what absolute validity
the knowledge won in all the special domains has, on what
presuppositions it rests, and what total conception of reality
it leads to. It is thus an analysis in the second power that is to be undertaken;
the work is not to be done over, but it is to be subjected, if possible, to
a reduction. And it is next a new synthesis that is to be undertaken,
and which is to embrace all the special syntheses. It is a task
that presents itself of necessity, and which cannot be put
off. But it is bound up with quite particular difficulties,
which philosophers have often taken very lightly.

We have remarked above that in synthesis it is especially the subjective
side in knowledge that plays a role. The unity and coherence
we seek in the phenomena is in its form determined by our spiritual
organization. Synthesis will therefore easily become more than an inverted
repetition of what already emerged in the analytic investigation;
there comes to be a surplus, which derives from the subject. But
how is this plus to be checked, so that it shall not falsify
the given? As long as we move only within a limited sphere,
% --- p. 59 ---
\setcounter{footnote}{0}%
\opage{59}there is here no insuperable difficulty. The coherence that we
posit is then itself a link in a larger coherence, within which it
stands in definite relations; and here it will then be possible to test whether
our conception of unity is justified or not. But when synthesis
is to embrace all that is given, when we ask for an explanation of the world,
such a check is impossible. We have nothing with which we could compare
our total conception of the universe, or to which we could bring
it into relation. The highest synthesis leads us to a limiting concept, to
ideas that lie at the limit of thinking and knowledge, beyond which
we cannot come by any means whatever. The activity of knowing
is essentially conditioned by the conscious being's relation to
the environment that asserts itself for it; but at the summit
where we here stand, the objective side of the relation drops away, and only the
subjective side remains. No testing by experience, no
verification is any longer possible.

Synthesis here necessarily becomes hypothesis. Hypothesis has
its fully justified place in science as an attempt at explanation, where
what is at hand is still insufficient to ground a real
knowledge. Inquiry then seeks to gather the given under a point of view
that agrees as much as possible with the facts. It
continues, only with dotted lines, so to speak, the lines it has begun
to draw within the verified knowledge. Even of
some of the special hypotheses (e.g. the nebular hypothesis, the Darwinian
hypothesis) it holds that they can be proved only with difficulty, perhaps not at all. But the possibilities of explanation that
can come into question in such special
domains are nevertheless only few, and the room for play will
constantly narrow more and more. One might therefore call these hypotheses
real, in contrast to those which seek to set up principles
for the explanation of reality as a whole. These will always retain
their ideal and subjective character. Therefore the thinker's personality
will here acquire an essential influence; the objective world-view
will receive an imprint of the subjective view of life.

Speculative philosophy has, however, hitherto in no way been willing to concede
this. It has always in its highest thought immediately seen the
highest reality, and it rejects with indignation any demand
for closer proof; it is precisely to be the distinguishing mark of speculative knowledge
that it has its truth in itself, needs no external
criteria of its validity. But it does not befit philosophy,
which above all should be self-conscious knowledge, and whose greatest
figure made ``know thyself!'' his motto, to conceal the real
coherence from itself.

In its conception of the method of philosophy Swedish speculation
shows, however, greater circumspection than the German. To be sure, the latter
also acknowledges the analysis of the given as a presupposition of synthetic
% --- p. 60 ---
\setcounter{footnote}{0}%
\opage{60}speculation; but it soon forgets this foundation and seizes a principle
from which it can independently construct itself and everything
real. ``Nothing is \textit{a posteriori}, everything is \textit{a priori}, otherwise pitiful half-measures!''
exclaimed the energetic Fichte, who was the founder of German speculation.
Without presuppositions, wholly from its own inner source, the
concepts were to well forth by an evolution embracing the whole content of existence.
Method and system were to be one, the movement of thought one with
the movement of the content itself.

But criticism\footnote[1]{\textit{J. F. Fries} in his critique of Schelling (``Reinhold, Fichte und Schelling''; Leipzig 1803). --- \textit{Trendelenburg} in his critique of Hegel (``Logische Untersuchungen'', 1840).} showed that all real determinations of content
were nevertheless taken over without further ado (and often even in distorted form) from
experience. It was thus only an illusion when thought imagined
that it gave birth to the content out of itself.

It is precisely one of the objections of the Boströmian philosophy to
German speculation that the latter lets the concrete develop out of the
abstract, or that it confuses ground of knowledge with real ground.
The abstract, says Prof. Nyblæus in his critique of Höijer (who indeed in
this respect joins the Germans), always presupposes the concrete.
Only the absolutely concrete (i.e. the absolute itself) is presuppositionless.
The concrete is real ground for the abstract; the abstract is for us,
on account of our limitation, ground of knowledge for the concrete,
in that by analysis (abstraction and reflection) we separate out the single determinations
from the concrete and examine them by themselves. But from the abstract
principles thus won we cannot immediately by synthesis
reach the concrete (the more concrete or content-richer determinations).
A concept, namely, can never be immediately developed out of another,
but [here we come to what is peculiar to the Boströmian theory of knowledge]
only out of the dimmer forms of self-consciousness, called
representation and feeling. The progress from the principle of science (the
simplest, most abstract determination) to what is conditioned by the principle
(``det principierade'') presupposes, therefore, that the object of science with
its whole concrete content is given to the thinker in the form of representation
or feeling, and the thinker must by continued analysis of this concrete
content (from which he had already won the principle) successively
develop the remaining, more concrete concepts that are to be combined with the principle
and grasped as its closer determinations. It will then appear
that there is a definite order that must be followed in the combination of the elements,
in that the more abstract determination must always in the concept precede
the more concrete as its formal ground. (``Den filosofiska
forskningen'', I, 2 p. 445—447).

There are in this presentation, in the nature of the case, two things to be examined
more closely: the given content and its analysis on the one hand, the
concepts won and their significance on the other hand.

% --- p. 61 ---
\setcounter{footnote}{0}%
\opage{61}For a method such as the one indicated to have any value, the first requirement
is that the content from which one proceeds be really given,
and given in definite forms. Before every obscurity in this respect
has been removed, the analysis cannot rightfully be concluded. Now the
immediate content that is given when philosophizing thought begins
its work is of a twofold kind. There is, first, the so-called
external world, a series of images and movements that unfold in time
and space, --- and there are, second, the subject's own inner states,
its thoughts, feelings and expressions of will, which succeed one another in time and
are given form only symbolically as images in space. But the philosophical
analysis does not here come upon raw material. The individual branches of science analyze
each its own circle of the content of the inner and of the outer world; the
results of this work the philosophical analysis must acknowledge (unless
in speculative arrogance it wishes to do it over), and it thus becomes
--- as we expressed ourselves above --- analysis to the second power, when it
attempts to use the results of the special analyses for a total analysis,
for the formation of a world-view. But the Boströmian philosophy
seems to hold that the situation is the same as in Plato's day, when the
philosopher was the only investigator, when the division of labor had not yet been
introduced in the domain of science. It designates the science of experience
with the derisive word ``empiry,'' which for it denotes the
counterpart of ``science.'' Empiry, says Boström, has to do with the sensuous,
speculation with the rational, the intelligible. Empiry gives only
the materials. An empirical philosophy has no more right to the name
of philosophy than superstition has to the name of religion. (``Grundlinier'' p. 9, 11).
% sv: blott falska eller empiriska orsaker --- causæ efficientes --- i stället för de sanna eller rationella
% sv: skenorsaker, hvilka den empiriska ytligheten anser för verkliga
% sv: högre syftning
The causes that empiry is able to point out are false, are only apparent causes\footnote[1]{``Grundlinier'' p. 22 (\dots{}. ``merely false or empirical causes --- \textit{causæ efficientes} --- instead of the true or rational ones''), p. 50 (\dots{}. ``apparent causes, which empirical superficiality regards as real''); cf. p. 55. --- \textit{Grubbe} directly sets ``empiricism'' in opposition to ``scientificity,'' which coincides with speculation and the ``higher aim.'' (II p. 323).};
only philosophy finds the true, rational causes (\textit{causæ finales}). --- Here, however, an essential intermediate determination is lacking. The polemic
against empiry is justified when by empiry is meant a loose
and random gathering of facts; Bacon already polemicized against what
he called \textit{experientia vaga}. But experience in the sense of exact
experimental determination of law and connection is something wholly different
from that loose gathering, which indeed falls outside of science.
Speculative philosophy makes things far too easy for itself by overlooking this
essential difference. If it would descend to earth and recognize
its community with earthly inquiry (it often seems as if it
--- like the Indian gods transformed into men --- had difficulty
standing firm on the earth), then it would acknowledge it as
a great good that the science of experience arranges the material
as it does. The given content with which philosophy (as speculation)
% --- p. 62 ---
\setcounter{footnote}{0}%
\opage{62}has to do is thus that which is given in the science of experience --- and it
offers a very different resistance to speculative treatment than
\textit{experientia vaga}, where there are no definite laws that could be broken.

What has been remarked here holds, to be sure, in a special degree of external nature,
but it finds full application in the psychological domain as well.
Empirical psychology is still so far behind that one would
think it must be of the greatest importance for philosophy to devote
as many forces as possible to it, in order to secure a solid point of departure
for its speculations. The whole character of modern philosophy
is conditioned by the consciousness of the necessity of such a psychological
foundation, and this is also acknowledged in the Boströmian method,
as Prof. Nyblæus has presented it. But the difficulty
that so easily leads the psychologist astray, namely the confusion of
the properly psychological facts with their explanation or
with the value one sets upon them, --- this difficulty the
Boströmian psychology succumbs to in a high degree. It presupposes that
in the immediate feeling there lies concealed an original content of ideas, which
it is merely a matter of bringing to the light of day. It distinguishes between a
lower and a higher side in the nature of man, of which the former is supposed to be
directed toward the sensuous, the latter toward the non-sensuous.
The difference between these two sides is qualitative, i.e. concerns not only the degree, but the kind and
content. Man has thus one content, one world in his lower feelings,
another in his higher feelings. (Nyblæus, I, 1 p. 39, 172, 192
etc.). But here psychology lacks exact investigation. What
does it mean that a feeling has a content? Feeling is in itself only
feeling, a purely subjective state; what has a content are the representations
and thoughts that evoke it, or which it evokes.
But whence do these representations in turn come, and whence do they get their
different content? Here there will be room for many kinds of more special
investigations, not only psychological in the narrower sense, but
historical as well. An immediately given content of consciousness is never
an ultimate that one can without further ado take as a foundation; it has its historical
origin and development, and only when one knows these does one know
it. It is thus very precarious to draw metaphysical conclusions
from such a content of consciousness --- and especially when aesthetic and ethical
valuations play a part. Concepts such as ``higher'' and
``lower'' are precisely only expressions of the value the phenomena concerned
may have from an ethical point of view, but ought to play no role
in psychology as a purely theoretical discipline. Such a division would
correspond to an anthropologist's dividing the races of men or a botanist
the plants into beautiful and ugly. And one can in any case
not infer from these valuations that there is a qualitative difference
between the phenomena: as beautiful and ugly plants may be very
closely related, so psychical phenomena, purely psychologically,
% --- p. 63 ---
\setcounter{footnote}{0}%
\opage{63}may stand very close to one another and yet, ethically considered, belong to very different
classes.

If the analysis were to have real philosophical significance, the
content of experience would first have to be spread out before our eyes with more or
less fullness, and it would have to be shown how the individual moments
can be developed from it. But we have already seen examples of
how the concepts in Boström undergo a very questionable metamorphosis
when they are taken out of the circle of experience. With the remark
that movement and change are not essential determinations in the concept of
life, he thinks he has established the justification for using this
as an expression for the Absolute. A study of recent biology shows
now, on the contrary, that life always presupposes a duality: namely, besides
the living organism, an environment (le milieu), a system of external
conditions, conditions de l'existence (Cuvier). And something similar holds
of the very concept of consciousness, as several philosophers (most
energetically Fichte) have shown.\footnote[1]{Boströmianism conceives the individual personality as an organ of a greater totality, a higher organism. But if the concept of personality were now applied to the Absolute, then this too would have to be the organ of something higher --- which becomes meaningless.}

It would probably appear, by a thoroughgoing analysis of the content of our consciousness,
that there is not a single determination that is not in its
essence relative and therefore cannot be applied to the Absolute\footnote[2]{And it does not help if one, as Stuart Mill proposed in his (on this point) unfortunate polemic against Hamilton, were to put the adjective ``absolute'' in place of the abstraction ``the absolute.'' For concepts such as absolute power, wisdom etc. suffer from exactly the same difficulties as the very concept ``the Absolute''; only the adjective conceals them better.}. It
was the insight into this, and not the assumption of a ``thing in itself'' (Nyblæus,
I, 2 p. 253), that led Kant to deny the possibility of a metaphysics
as knowledge of the essence of things. This insight was still alive
in Fichte and was contending in him with the speculative impetus. Not until
Schelling and Hegel did men set themselves wholly above these scruples.
Later Hamilton, Bain, Spencer and Wundt in particular have demonstrated
the thoroughgoing logical and psychological validity of the law of relativity.
Kant himself thought, to be sure, that there were certain privileged psychological
facts that were to be capable of being emancipated and of emancipating us from
this law. The feeling of duty, moral responsibility --- a moment of what
in Boströmianism is called ``higher experience'' --- opened for him
the entrance to the world of the Absolute. In a similar way Hamilton and
several Neo-Kantians. But here it is precisely a matter of not letting oneself
be led astray by valuations and of not without further ado transcribing
them into theory. This is the fertile source of dogmas, a refuge for the \textit{ignava ratio}. However great or however small the value a fact has for us,
we can nevertheless learn from it not the least about its explanation or
about whether or not it can be explained by natural causes.
When the philosophical discussion enters the domain of
% --- p. 64 ---
\setcounter{footnote}{0}%
\opage{64}valuations, and uses the designations ``higher'' and ``lower'' as arguments,
it is on a precarious slope. One then polemicizes,
like Bishop Mynster among us, against Locke as a kinsman of ``all bad
and contemptible philosophemes,'' or one addresses one's auditorium, as Victor
Cousin did his in 1815: ``Messieurs, vous aimez ardemment la patrie. Si vous
voulez la sauver, embrassez nos \emph{belles} doctrines.'' Thinking ceases,
and agitation begins. --- Valuations have their full
justification within philosophy too; it is the task of ethics and aesthetics
to investigate and ground them, but they cannot be used
to help with the difficulties of psychology and metaphysics (not to speak of
physics).

But if it now succeeds --- to come to the second main point
of the Boströmian method --- in separating the series of fundamental concepts
from the empirical content and bringing them into inner harmony
with one another, what validity can then be ascribed to these concepts?
--- When the anatomist has isolated and prepared various elements
in the organism, he does not suppose thereby to have arrived at a new and higher
organism. By considering the individual anatomical elements by themselves
he better understands their position and role within the organism and makes
it easier for the physiologist to construct for himself the inner life of the organism.
But should the philosopher, who seeks to isolate the fundamental determinations
of human consciousness, be differently placed? By his
analysis he can gain deeper insight into the inner coherence of the life of thought, its development
and its range; but is there anything whatever that entitles
him to regard the abstract fundamental concepts as expressions of
a real, higher and more comprehensive life of thought? If he does so, he takes
here a leap that he himself does not notice or at any rate
does not explain. It is a mystical wave that seizes him at this
point and carries him away. Abstraction is abstraction, and can
never become anything else. That hypostatizing of the abstract ideas must be due to
a power other than theoretical knowledge itself.

We miss at this point a definite indication of the way by
which knowledge is able to attain a positive concept of the Absolute
as the true reality. And this is all the more remarkable
as the Boströmian philosophy otherwise lays great weight on the immanence
of consciousness and polemicizes so eagerly against popular realism
because it goes behind consciousness. Boström really gives only a
mystical solution of this difficulty. It lies in the concept of consciousness,
he says, that it can sense only itself and its own
sensations. The presupposition for the human spirit's being able to sense
God is therefore that man's own consciousness is in a certain
respect (aliqva ex parte) identical with God's consciousness. (``Om
relig.'' etc., p. 164). But whence is it known that this presupposition
holds? That my consciousness is to a certain degree identical with the
% --- p. 65 ---
\setcounter{footnote}{0}%
\opage{65}consciousness of another being, I can know only when I can compare
both the two consciousnesses; but then it is presupposed that I know them
both beforehand. From this circle we could not escape.

In reality all speculative philosophy has a religious character. It
demands a giving up of reflection, so that consciousness can immediately
apprehend its unity with the Absolute. Therefore it asserts that, besides
discursive knowledge, there is a higher kind of knowledge that
grasps the Absolute, --- an intuitive knowledge, an intellectual intuition, or
however one may now designate it. And only thereby can it open a prospect
of a higher content of reality, a positive doctrine of the Absolute. But
this advantage is bought only by giving up what had been won by
the struggle against realism, since speculation itself now sins against
the law of subjectivity by going behind consciousness. For it comes
to the same thing whether one posits an absolute spirit or an absolute
matter; in both cases one commits the same surreptitious assumption in the theory of knowledge.
As Zeus boasts that he will draw the whole world up
to himself by a golden chain, fasten it to one of Olympus's peaks and
thereby set everything in suspension, so it has always been the task
of speculation to find an absolute point about which all existence
might turn. But our position is, after all, not the Olympian one.
The point to which we seek to fasten everything is always itself movable.

Boström himself concedes that consciousness, as we know
it, is always imperfect and relative. Human striving
aims at attaining as great clarity and freedom as possible; with clarity
and freedom the highest possible unity and harmony in
all spiritual functions is then bound up. (``Om relig.'', p. 152). But by negation
of our limitation and imperfection we form the concept of the
perfect: when it is relations that are negated, the negation is a perfection.
(``Om relig.'', p. 108). --- But if Boström here does not
wish to swear by Hegel's banner, it is impossible for him to attain a new positivity
by way of negation. By negation of the imperfect we in reality obtain
only the perfect in the sense of the not-imperfect. We
think the Absolute only in so far as we negate the relative. But
since there is no wholly negative thinking (such would be
one with not-thinking), one always calls the imagination to one's aid when
one wishes to think the Absolute. One represents to oneself more or less
distinctly an image, a symbol, which reflection immediately corrects
and dissolves, but only so that it may be posited by the imagination in an
even more sublime form, again dissolved --- and so on to infinity.
As Antaeus always drew new strength by contact with the earth,
so the imagination draws all the elements for the idea of the Absolute from
the world of experience; but thought consumes them by raising them beyond
the determinate and limited sphere within which they have their
positive significance.

% --- p. 66 ---
\setcounter{footnote}{0}%
\opage{66}The secret of speculation is that abstraction stands in the service
of idealization. Every time we form an ideal for ourselves, we disregard the
accidental and limited that clings to every given phenomenon; we
hold to certain properties and characteristic marks that have our full
assent, and think how a being would be when it merely
had these traits, and had them in the fullest measure and in mutual harmony.
This idealizing process can be clearly traced in the Boströmian
philosophy. The spiritual life, as experience shows it to us, possesses
always only a relative perfection; it struggles with oppositions
and contradictions; it lacks true unity and harmony. The
highest that we find here is a constant striving after reason and freedom.
``The human spirit possesses a certain measure of life and self-consciousness, beyond
which it is not able to develop in any form of life whatever;
however high a degree of perfection it may attain, there always remains
something more imperfectly apprehended.'' (``Om relig.'', p. 162 f.). Experience
thus shows us a scale of stages of development, whose conclusion is
unthinkable so long as the talk is of finite beings. (``Grundlinier'' §
56). But with a brisk jerk the imagination then puts the absolute superlative
in place of these comparatives; or rather: all the comparatives
are ordered in a system and seen as comprised in an eternal superlative. E.g.: man is a rational being, but not reason itself. There
are in his essence other elements than reason. Only the Godhead is
absolute reason, a sensing being without any limitation. (``Satser
om lag och lagstiftning'', p. 22 f.). We stand here again, then, at the proposition
that the imperfect is understood by the perfect, and we can now
rewrite it thus: Reality is understood by the ideal.

There is in man a need to think of the end of all
his endeavors as finished and complete. In thought and longing he occupies himself
with that which he still lacks and which he struggles to attain;
and besides, he must give it a certain degree of definiteness for himself
in order to be able to find means and ways to reach it. His present
condition is one of want and struggle; but the condition he has in view is
the perfectly harmonious life resting in itself. And now the relation
reverses itself for him in such a way that it looks as if the end were
an already existing reality, from which he is separated only by his
imperfection, so that he need only think this away in order to
have the idea of the former. The ideal existence has for him in a far deeper
sense reality than empirical imperfection, which he does not
deem worthy of the predicate of true reality.

That the very concept ``reality'' thus becomes a valuation
shows that it is a practical motive that underlies speculation.
There can be traced in the speculative systems an undercurrent of
feeling and enthusiasm, which often explains what the systems' own concepts
do not make intelligible; and the lasting value of these systems rests
% --- p. 67 ---
\setcounter{footnote}{0}%
\opage{67}not least on their being expressions of such unconscious tendencies in
the life of feeling. But this also implies that such systems cannot
lay claim to strict scientific validity --- let alone to
have the right to regard themselves as the expression of the true and highest
science.

Prof. Nyblæus, who strongly emphasizes the practical-ethical
tendency of Swedish speculation, has shown that the idea of the harmonious totality
is its proper fundamental thought. In Leopold, Thorild, Boëthius
and Höijer we see it already standing in the foreground, and in Boström it receives
its full development. It is an ethical idea of the very greatest significance.
The unhistorical conception and oppositional
tendency of the eighteenth century had the effect that for a long time this idea was misjudged and it was thought
possible to derive the ethical from the nature of the single individual. But
whichever side of the individual's nature one then holds to, one will always
be led further back and discover that the individual is understood only as a member of
the whole race from which it springs, whose heir it is, and whose
tasks it must work on in order to be able to develop its own essence. A
harmonious society of free personal beings may be said to be the
highest ethical ideal for human action. But what right
have we to regard this ideal as really existing? It is a
model that we project for our action; we have it before us, but it
rests only on an illusion when we at the same time place it behind us as an absolute
presupposition. And by such a hypostatization the ethical interest
will gain nothing, but only lose. It will gain nothing. For it gives
me no deeper insight into the essence of ethical action and of the ethical
forms of life that I, with Boström, think of society as
a higher, more comprehensive personality, in which I am a moment. Boström
does not even himself make use of this mystical presupposition in
his investigations. When he wishes, e.g., to ground his doctrine of hereditary
monarchy as the right constitution of the state, he relies on
grounds drawn from the demands of real social life, in that he in particular
seeks to show that the ruler must be raised above all private and
particular interests, in order to be able, wholly disinterested, to represent
the public and universal interests, --- and that the welfare of the state requires
absolute unity in the administration, which is not supposed to be attainable in a republic.
The republic --- he thinks --- corresponds to that age in the individual
human being when the passions struggle with one another, whereas monarchy
is the expression of the dominion of the rational will. (``Satser om
lag och lagstiftning'', 2. uppl., p. 33 ---35, 63—64). If this is
right, it can be proved only by psychological and historical grounds,
and such are in reality constantly appealed to. In a
similar way matters stand with the justification of penal right.
The Boströmian philosophy proves its humane character by polemicizing
against the theory of absolute retribution, in contrast to which it
% --- p. 68 ---
\setcounter{footnote}{0}%
\opage{68}maintains: that an evil, as punishment is, cannot be the final end, but
can only be defended as a necessary condition for a greater good. From
this proposition Boström argues in his sharp polemic against
the theological doctrine of eternal damnation. And in it there lies
also all that is needed for a justification of the state's penal right,
without its becoming necessary to conceive crime as an injury
to the state's mystical personality. In any case it is always precarious
to build on so hypothetical and speculative a foundation in so
important a matter. The conception of the state as a harmonious whole, a total
``personality,'' indicates only the goal that is to be reached; the means of reaching
this goal can be taught us only by progressive historical experience, ---
just as it is also in experience that they must stand their test.\footnote[1]{Cf. \textit{Hans Forssell:} ``Studier och kritiker'' (Stockholm 1875) p. 174---178.}

The ethical interest, on the other hand, will precisely be weakened by
its ideal's being posited as an eternal being. Why should I work,
when everything I strive after already is? That which is obscure,
imperfect and disharmonious in me is, according to Boströmianism,
mere phenomenon, due only to my finite conception; my true essence
is an eternal, harmonious reality. But once I have grasped
this, i.e. once I have become a Boströmian, I can calmly let that
phenomenon be what it is, in order to live wholly in the intuition of my true
self. And if one would say that the opposition between the clear and the
obscure, between my essence and my phenomenon, will not leave me at rest in this,
and that I must constantly work theoretically and practically to remove
the obscurity, then it is not easy to see what that hypostatizing of
the ideal helps me, when it cannot at all be attained and realized
by me. Even in its highest possible form of life the finite being, according to
Boströmianism, is separated by an insuperable gulf from the absolutely
ideal existence. But what advantage then has the hyperidealist
over us poor empiricists, who can indeed form ideal models,
but cannot ascribe to them any other validity than precisely
that of being rules and guiding stars for our endeavors? Both
parties concede that the opposition between ideal and reality is
insuperable for us; but the hyperidealist nevertheless thinks he can
attain a knowledge of the ideal as the only true reality, whereas
the empiricist finds in that insuperability a new proof of the law of
relativity and an expression of the fact that existence for us is striving, movement
and development. The abiding opposition between ideal and
reality contains a sting that we must constantly feel anew. It
is of no use for Boström to console us with the fact that we do after all reach in the end
the goal, so far as it can be reached by finite beings, and that we therefore
can no longer have any wish. For psychology teaches us
that so long as we are able to apprehend the difference between ideal and
% --- p. 69 ---
\setcounter{footnote}{0}%
\opage{69}reality, so long we also regard reality as imperfect. Therefore
conclusion, consummation, absolute satisfaction stand for all striving
natures as equivalent to annihilation. Lessing's saying is
well known: that he would prefer the eternal striving after truth to
the possession of absolute truth, and that eternal bliss would
be one with eternal boredom. When Duroc said one day to
Napoleon that people considered him so ambitious that he would gladly, if it were
possible, take the place of God the Father, he answered: ``\textit{Ah! je n'en voudrais \emph{pas; c'est un cul-de-sac}}.'' This is not merely an expression
of ambitious insatiability, but aptly designates the fundamental
law of our existence: every result can for us be only the foundation of a new beginning.
However far we move the boundary, the same will repeat itself.

The law of relativity teaches us that our whole conception of the world is subjectively
conditioned, holds only from our standpoint and is not the only
possible one. We then try to raise ourselves above this limitation, and believe
by abstraction from our imperfection, by the ``negation of the negation,''
to attain an absolute. But every new concept that we believe by this way to
win always proves to be subject to the same limitation and relativity.
It is the old story of the house-sprite that always moves along too.
Speculative philosophy believes that our limitation is one thing, our true
essence another, and that one can separate the former and hold to the latter.
But we have precisely our strength in limitation; the law of relativity is
not merely an expression of our barriers, but also of the positive means
at our command. Without barrier and resistance no thinking and acting.
Relativity is fruitful, striving and developing, whereas
absolutism leads us into a theoretical and practical cul-de-sac. The
Absolute is for us only a \textit{negative} concept: we think it in so far as
we think our own limitation. But as such it is of no small
significance. It leads us constantly beyond the immediately given, and is
therefore a principle of constantly renewed theoretical and practical striving; we see
ourselves as ``one among many,'' as occupying a single, limited position
in the great manifold of existence. The validity of our theoretical
insight and of our practical ideals we could not establish by
comparing them with an absolute (whether this be absolute reason
or absolute matter), but only by their proving themselves, under the ordeal by fire of continually progressing
experience, to lead to greater coherence and
agreement among our representations and to greater harmony and
richer development of human life as a whole. But a conclusion
to this process is for us unthinkable, is one with absolute
death, the absolute nothing.

Swedish speculation is, as we have seen, in its method not so
bold as the German, does not set itself such great tasks as the latter does.
Both schools are uncritical in their defective and unsuccessful grounding
of their absolute principle; but they part ways in the further
% --- p. 70 ---
\setcounter{footnote}{0}%
\opage{70}execution, in that Boström thinks he can help himself with analysis,
whereas Schelling and Hegel demand a synthetic construction. This
difference in method is connected with an equally essential
difference in the world-view. According to Schelling and Hegel the speculative
construction is not merely an expression of the thinker's subjective reflection,
but coincides with the Absolute's own self-development.
The being of the Absolute is an eternal process. But Boström rejects this
as empiricism. That we conceive existence as development is supposed to
have its ground only in our finitude and limitation; we are not able to
grasp the eternal content of ideas in a clear and true way.
Change and development do not concern the content itself, but our conception of it.
The proper development therefore consists in the transition of our knowledge from
obscurity to clarity, whereby the eternal, absolutely original content
is more and more illuminated. But by this process of clarification the Absolute
is wholly untouched; it is in itself always complete, whereas the sensuous
never in reality becomes complete. (``Grundlinier'' § 1, 6, 35, 38).

It is not only against German speculation that Boströmianism
here declares itself, in that it accuses it of making the Absolute finite;
it is also against Christian theology. The principal Christian dogmas
(in particular the dogma of creation and the dogma of atonement) rest on the
presupposition that God can act (produce, create) and suffer (be offended and
reconciled), and in this they sin against the right and pure concept
of an absolute being. Speculative pantheism is --- as Prof.
Nyblæus has shown in earlier treatises --- only a direct consequence
of the Christian dogmas. One cannot, he says, ascribe to
God a producing or creating activity, cannot let the sensuous
world be his immediate work, without his essence's being confounded
with it. For what appears as the result of an activity
must be contained potentially in the acting force. The activity
is only a transposition from possibility to actuality of an original
content. He who lets God create or produce the world
therefore makes the world one with him. A true distinction between the
infinite and the finite, the spiritual and the sensuous, becomes only
possible when one assumes that God and the ideal world are eternally complete
and real, and that it is only in and for our sensuous apprehension that
a process of development takes place. An infinite development, moreover, conflicts
with the very concept of development; for when the goal lies infinitely far
away, it can never be reached, and one thus lacks a standard by which to decide
whether a development from lower to higher stages is really taking place.
An infinite force could just as little be thought: it would at once
be at the goal, its effects would coincide with itself\footnote[1]{\textit{Axel Nyblæus:} ``Trenne religionsfilosofiska uppsatser'' (Lund 1872) p. 116 f, 126 f, 141, 145. --- ``Om politisk frihet och sjelfstyrelse'' (Lund 1869) p. 39, (``Trenne
\opage{71}filosofiska uppsatser'', Lund 1878, p. 43). --- ``Den filosofiska forskningen i Sverige'', I, 2 p. 275, 279, 427.}.

% --- p. 71 ---
\setcounter{footnote}{0}%
\opage{71}It cannot surely be denied that if one wishes to think the concept of the Absolute
consistently and purely, no development in time can be ascribed to it,
because there can be no \textit{plus ultra} for it. Boströmianism
here shows a scruple that it does not show at other
points; for --- as remarked earlier --- the same difficulties
as with this concept confront the philosophy of the Absolute with the concepts
``life'' and ``self-consciousness.'' Plotinus saw this: Since thinking
is an activity, it cannot be ascribed to the highest principle. (``Enneades''
III. 9, 3). If Boström is more consistent than Hegel, then Plotinus is
more consistent than Boström, because he saw that \textit{none} of our concepts
suffices for the demand to think the Absolute.

But it is a question whether it is justified to ascribe to Hegel
the doctrine that the Absolute develops in time. Its development is according
to Hegel purely logical, is only an expression of the fact that the one moment in its
concept leads with inner necessity to the others, so that at whatever
point one wishes to begin, one will be led through the whole system of
its determinations. Hegel even flatly denies a \textit{real} process of development
in nature. ``Nature,'' he says, ``is to be regarded as a \textit{system
of stages,} of which the one necessarily proceeds from the other and
is the nearest truth of that from which it results, yet not in such a way
that the one was naturally \textit{produced} (erzeugt) of \textit{the other, but
in the inner idea,} which constitutes the ground of nature''\footnote[1]{\textit{Hegel:} ``Encyclopädie der philosophischen Wissenschaften'', 4te Aufl. (Berlin 1845) p. 204 (§ 249).}. When the germ develops
into a plant, Hegel therefore derives this from the logical
concept of the germ itself: the germ is the contradiction that it is only in itself and yet
is not supposed to be so\footnote[2]{\textit{Hegel:} ``Geschichte der Philosophie'', 2te Aufl. (Berlin 1840) I. p. 35.}. The transition is thus \textit{properly} purely logical, and
the difficulty about time Hegel can so far just as well push
aside as Boström does when he ascribes to the Absolute thinking
and self-consciousness. Sibbern rightly missed in Hegel development
in the sense of a real \textit{historical} process\footnote[3]{\textit{Sibbern:} ``Speculativ Kosmogonie'', (Copenhagen 1846) p. 7.}.

Even in the domain where Boströmianism posits a development,
this has no historical character. The content is indeed eternal and
complete; what goes on during development is only, so to speak,
an uncovering or peeling off of enveloping husks. Boströmianism
will not concede that anything really \textit{happens}, i.e. becomes something
that was not before. Therefore, e.g., Prof. Nyblæus reproaches Locke with
assuming a real origin of our cognitions instead of a
mere transition from obscurity to clarity of an original content lying
in us. (I. 1, p. 3). What comes to be must always have been,
only obscured. This points to a concept of development that is now
% --- p. 72 ---
\setcounter{footnote}{0}%
\opage{72}abandoned by all, with the exception of speculative philosophers. It seems to
be something that for many is bound up with peculiar difficulties,
to concede that as the result of development something can come forth
that in no sense existed beforehand. When in spring the child sees the leaves
come forth, it thinks that they have lain inside the branches, and that it is now
opened for them. Its reasoning is that in order for them to be able to come
out, they must first have been brought in there. Expressed in a more learned way
this means that what is actual must already have existed potentially.
In this way everything gets an eternal preexistence. Thus
Leibnitz assumed that reason was originally laid into the germs that develop
into rational beings, so that development (développement) presupposes involution
\textit{(enveloppement}\footnote[1]{``\textit{J'aimerois me passer des miracles dans la génération de l'homme, comme dans celle des autres animaux; et cela se pourra expliquer, en concevant que dans ce grand nombre d'âmes et d'animaux, ou du moins de corps vivans qui sont dans les semences, ces âmes seules qui sont destinées à parvenir un jour à la nature humaine,} \emph{\textit{enveloppent}} \textit{la Raison qui y paraîtra un jour}''. --- ``Théodicée'' § 397; cf. ``Monadologie'' § 73---75.}. It is this theory that has been called the box hypothesis,
because it imagines things as hidden inside one another and development as
consisting in taking them apart. The later progress of the sciences
has led to abandoning this concept of development, which
abolishes all qualitative change. Kant was the first to set up
another type of development, in that in his ``Allgemeine Naturgeschichte
und Theorie des Himmels'' (1755) he sets up the well-known, still current
hypothesis of the origin of the solar system from a chaotic nebular mass
by progressive articulation. Shortly afterwards (1759) the anatomist
C. F. Wolff, the founder of scientific embryology, showed that the fetus arises
successively from a uniform rudiment, and set up the theory of epigenesis,
of real new formation, in opposition to the so-called theory of evolution,
which taught a mere uncovering of parts already formed. This theory
has --- despite Haller and Bonnet --- broken through and has received its full
confirmation\footnote[2]{Cf. \textit{Kölliker:} ``Entwicklungsgeschichte des Menschen und der höheren Thiere'' (Leipzig 1861) p. 4.}. He who wishes to form his concepts according to experience has
no other way out than to follow the track thus struck out. Herbert
Spencer has shown how the concept of development in its main features is
one and the same in all the different domains, and that its essential
characteristic marks are transition from a uniform to a multiform (differentiated)
state, from less to greater coherence, from indefiniteness to definiteness\footnote[3]{See my book: ``Den engelske Filosofi i vor Tid'' (Copenhagen 1874) p. 136---138, 160---167.}. The spiritual and the material world, as we
know them by the way of experience, show the same type in their process of development.

In face of the wealth of features and events that development in the world of experience
shows us, it looks poor when Boströmianism
has nothing else to tell us than that the whole is a transition from
% --- p. 73 ---
\setcounter{footnote}{0}%
\opage{73}obscurity to clarity. It is a childlike speculation's attempt to escape
from the paradox --- a paradox that runs through all our knowledge
from first to last ---, that something really comes to be, thoughts and cognitions
as well as plants, animals and worlds.

German speculation shows itself here to stand nearer to experience
than the Swedish. It is so vividly convinced of the significance of development
that it --- quite inconsistently, to be sure --- transfers it to the
Absolute. The Christian doctrine of a suffering God points in the same
direction, expresses that it is once and for all a fundamental law that perfection
is reached only through development, is not eternal. Seen from the practical standpoint,
this conviction is certainly of far greater significance (despite
its metaphysical inconsistency) than the pure world of ideas of Plato and Boström,
untouched by all becoming, all strife and suffering. We ourselves are striving
and suffering beings, and it avails us but little to know that there is
such an eternal ideal world, when it is always raised above the course on
which we move. The practical motive from which Swedish speculation
has sprung has here shot past the mark and works against itself.
The world-view that rests on the fundamental law of relativity,
on the other hand, is able to keep the idealizing tendency in the human
spirit, this source of all poetry and free speculation, within the proper
bounds, whereby the one becomes able to work in the service of real ends.

\end{document}
